\documentclass[a4paper, 11pt]{book}
\usepackage{anyfontsize}
\usepackage[CJKchecksingle, AutoFakeBold]{xeCJK}
\usepackage{CJKnumb}

\usepackage{setspace}
\onehalfspacing
\renewcommand{\arraystretch}{1}

\setlength{\parindent}{2em}

\usepackage[intlimits, leqno]{amsmath}
\usepackage{amssymb}
\usepackage[centercolon]{mathtools}
\usepackage{siunitx}

\usepackage{fontspec}
\usepackage{unicode-math}

\setmainfont[
    Extension=.otf,
    UprightFont=*-Regular,
    BoldFont=*-Bold,
    ItalicFont=*-Italic,
    BoldItalicFont=*-BoldItalic
]{STIXTwoText}

\setmathfont[
    Extension=.otf,
    BoldFont=*
]{STIXTwoMath-Regular}

\let\exerciseQuestion\Question
\AtBeginDocument{%
    \let\mathQuestion\Question
    \let\Question\exerciseQuestion
}

\setsansfont{TeX Gyre Heros}
\setmonofont{Courier New}

\usepackage{esint}

\DeclareSymbolFont{CMlargesymbols}{OMX}{cmex}{m}{n}
\let\sumop\relax\let\prodop\relax
\DeclareMathSymbol{\sumop}{\mathop}{CMlargesymbols}{"50}
\DeclareMathSymbol{\prodop}{\mathop}{CMlargesymbols}{"51}

\usepackage{euscript}

\newcommand{\bm}{\symbfit}
\newcommand{\smalltriangle}{\mathbin{\vcenter{\hbox{\scalebox{0.65}{$\triangle$}}}}}

\setCJKmainfont[
    ItalicFont=FZKai-Z03S
]{FZShuSong-Z01S}

\setCJKsansfont[
    BoldFont=* Medium,
    ItalicFont=FZKai-Z03S
]{Noto Sans CJK SC}

\setCJKmonofont[
    ItalicFont=FZKai-Z03S
]{FZShuSong-Z01S}

\setCJKfamilyfont{kai}{FZKai-Z03S}

\setCJKfamilyfont{song}[
    ItalicFont=FZKai-Z03S
]{FZShuSong-Z01S}

\setCJKfamilyfont{fangsong}[
    ItalicFont=FZKai-Z03S
]{FZFangSong-Z02S}

\setCJKfamilyfont{hei}{Noto Sans CJK SC Medium}
\setCJKfamilyfont{hei2}[
    BoldFont=* Bold
]{Noto Sans CJK SC}

\setCJKfamilyfont{black}[
    BoldFont=* Black
]{Noto Sans CJK SC}

\setCJKfamilyfont{sectionfont}[
    BoldFont=* Black
]{Noto Sans CJK SC}

\setCJKfamilyfont{pffont}[
    BoldFont=* Medium
]{Noto Sans CJK SC}

\setCJKfamilyfont{emfont}[
    BoldFont=* Medium
]{Noto Sans CJK SC}

\defaultfontfeatures{Ligatures=TeX} 
\XeTeXlinebreaklocale "zh"
\XeTeXlinebreakskip=0pt plus 1pt minus 0.1pt

\newcommand\kaishu{\CJKfamily{kai}}
\newcommand\songti{\CJKfamily{song}}
\newcommand\heiti{\CJKfamily{hei}}
\newcommand\thmheiti{\CJKfamily{hei2}}
\newcommand\fangsong{\CJKfamily{fangsong}}
\renewcommand{\em}{\bfseries\CJKfamily{emfont}}


\usepackage[unicode, colorlinks, bookmarksnumbered]{hyperref}
\hypersetup{
	linkcolor=blue,
	citecolor=red,
	urlcolor=teal,
	pdfauthor={王畅},
	pdftitle={ICS 问题求解（2021 秋）}
}

\usepackage[iso, english]{isodate}
\usepackage{geometry}
\geometry{
	paper=a4paper,
	top=3cm,
	inner=2.54cm,
	outer=2.54cm,
	bottom=3cm,
	headheight=6ex,
	headsep=6ex,
}

\usepackage[svgnames, dvipsnames]{xcolor}

\usepackage{paralist}
\usepackage{commands}
\usepackage[font=small, labelfont=bf]{caption}
\usepackage{graphicx}
\usepackage{float}
\usepackage{longtable}
\usepackage{multirow}
\usepackage{tikz}
\usetikzlibrary{shapes.symbols, backgrounds, matrix, calc, arrows, math, arrows.meta, shapes.gates.logic.US, shapes.gates.logic.IEC, decorations.pathreplacing}
\usepackage{minted}
\renewcommand\fcolorbox[4][]{\textcolor[rgb]{0.63,0.63,0.00}{\strut#4}}

\newcommand\MemoryLayout[2]{
	\begin{tikzpicture}[scale=0.3]
		% \draw[thick] (0,0) -- ++(0,3) node[above] {$0$};
		\foreach \pt/\col/\lab [remember=\pt as \tp (initially 0)] in {#1} {
			\pgfmathparse{\pt-1}
			\foreach \a in {\tp, ..., \pgfmathresult} {
				\draw[fill=\col] (-\a,0) rectangle ++(-1,2);
			}
			% \draw[thick] (-\pt,0) -- ++(0,3) node[above] {$\pt$};
			\if\lab\relax\relax\else
				\draw[thick, decorate, decoration={brace, amplitude=4mm}] (-\tp,-0.2) -- node[below=4mm]{\lab} (-\pt,-0.2);
			\fi
    	}
		#2
  	\end{tikzpicture}
}

\numberwithin{equation}{chapter}
\renewcommand{\theequation}{\thechapter.\arabic{equation}}

\usepackage{ulem}
\normalem

\usepackage{datetime}

\usepackage{xstring}
\usepackage{fancyhdr}
\newcommand\sectionname{\arabic{chapter}.\arabic{section}}
\pagestyle{fancy}
\renewcommand{\chaptermark}[1]{\markboth{
	\CJKnumber{\thechapter}\quad #1
	}{}}
\renewcommand{\sectionmark}[1]{\markright{\arabic{chapter} \quad #1}}
\fancyhf{}
\fancyhead[EC]{\CJKfamily{hei2}\footnotesize{\leftmark}\vspace{1mm}}
\fancyhead[OC]{\CJKfamily{hei2}\footnotesize{\leftmark}\vspace{1mm}}
\fancyhead[LE, RO]{{\footnotesize \thepage}\vspace{1mm}}
\fancyhead[RE, LO]{}
\renewcommand{\headrulewidth}{0.5pt}
\renewcommand{\footrulewidth}{0pt}
\addtolength{\headheight}{0.5pt}
\renewcommand{\headrulewidth}{0pt}
\renewcommand{\footrulewidth}{0pt}
\addtolength{\headheight}{0.5pt}

\usepackage[many]{tcolorbox}

\newtcolorbox{summary}{
	breakable,
	enhanced,
	width=\textwidth,
	colback=white, colbacktitle=white,
	colframe=gray!50, boxrule=0.2mm,
	coltitle=black,
	fonttitle=\bfseries\sffamily,
	attach boxed title to top left={yshift=-\tcboxedtitleheight/2, xshift=\tcboxedtitlewidth/4},
	boxed title style={boxrule=0pt, colframe=white},
	after skip=1cm,
	top=3mm,
	bottom=3mm,
	title={要点}
}

\usepackage[calcwidth, explicit, nobottomtitles, newparttoc, indentafter]{titlesec}
\usepackage{titletoc}
\usepackage[titles]{tocloft}
\cftsetpnumwidth{1.75em}
\providecommand{\dmchapterttl}[1]{\IfSubStr{ABCDEFGHIJKLMNOPQRSTUVWXYZ}{#1}{附录 #1	}{\CJKnumber{#1}}}

\newlength{\BoxTtlwidth}

\newcommand{\MakeChapBox}[2]{%
	\settowidth{\BoxTtlwidth}{\dmchapterttl{#1}}
	\begin{tcolorbox}[
		enhanced jigsaw,
		skin=bicolor,
		frame engine=path,
		sharp corners=all,
		width=0.9\textwidth,
		top=4mm, bottom=4mm,
		sidebyside,
		frame hidden,
		boxrule=0mm,
		lefthand width=\BoxTtlwidth,
		colupper=white,
		colback=gray!80,
		colbacklower=gray!10,
		sidebyside align=center,
		halign=center]
		\dmchapterttl{#1}
		\tcblower
		#2
	\end{tcolorbox}%
}

\newcommand{\MakeChapBoxSingle}[1]{%
	\begin{tcolorbox}[
		enhanced,
		width=0.7\textwidth,
		sharp corners=all,
		top=4mm, bottom=4mm,
		frame hidden,
		boxrule=0mm,
		colback=gray!10,
		halign=center]
		#1
	\end{tcolorbox}
}

\newtcbox{\MakeSectBox}{
	enhanced,
	arc=0pt, outer arc=0pt,
	before skip=0pt, after skip=0.4em, left skip=0pt, right skip=0pt,
	top=0pt, left=0pt, right=0pt, bottom=1.5mm,
	sharp corners=all,
	valign=bottom,
	colback=white,
	colframe=white,
	boxsep=0pt, leftrule=0pt, rightrule=0pt, toprule=0pt, bottomrule=0pt,
	overlay={ \draw[line width=1pt] (interior.south west) -- (interior.south east); }
}

\newcommand*\chapterlabel{}
\titleformat{\chapter}
  	{\gdef\chapterlabel{}\sffamily\CJKfamily{black}\bfseries\Huge}
  	{\gdef\chapterlabel{\thechapter\ }}{0pt}
  	{\begin{tikzpicture}[remember picture, overlay]
		\node[yshift=-3cm] at (current page.north west)
	  		{\begin{tikzpicture}[remember picture, overlay]
				\draw[fill=titlepagecolor!20] (0,0) rectangle (\paperwidth, 3cm);
				\node[anchor=east, xshift=0.9\paperwidth, rectangle,  rounded corners=20pt, inner sep=11pt, fill=titlepagecolor] {\color{white} {\fontsize{48}{60}\selectfont\chapterlabel}#1};
			\end{tikzpicture}};
   	\end{tikzpicture}
}
\titlespacing*{name=\chapter, numberless}
	{1pc}{*4}{1em}

\titleformat{name=\section}
	{\filleft\normalfont\sffamily\bfseries\CJKfamily{sectionfont}}
	{}
	{0mm}
	{ \settowidth{\BoxTtlwidth}{\Huge \thesection \hspace{0.7em} \Large #1}
		\ifdim \BoxTtlwidth < \textwidth
			\MakeSectBox{\Huge \thesection \hspace{0.7em} \Large #1}\vskip-18pt%
		\else
			\Huge \underline{\thesection} \hspace{0.7em} \Large #1%
		\fi%
	}
	[]

\titleformat{name=\section, numberless}
	{\filleft\normalfont\sffamily\bfseries\CJKfamily{sectionfont}}
	{}
	{0mm}
	{ \settowidth{\BoxTtlwidth}{\Large #1}
		\ifdim \BoxTtlwidth < \textwidth
			\MakeSectBox{\Large #1}\vskip-18pt%
		\else
			\Large #1
		\fi
	}
	[]
	
\titlespacing*{name=\section}
	{1pc}{*1.3}{*1}

\usepackage{zhnumber}
\usepackage[answerdelayed, lastexercise]{exercise}
\renewcommand{\ExerciseListName}{}
\renewcommand{\AnswerListName}{}
\renewcommand{\ExerciseHeaderNB}{\theExercise}
\renewcommand{\QuestionNB}{(\arabic{Question})}
\renewcommand{\ExerciseHeaderOrigin}{\ ({\kaishu\ExerciseOrigin})}
\renewcommand{\ExerciseListHeader}{\ExerciseHeaderDifficulty\textbf{\ExerciseHeaderNB}\textbf{.}\ExerciseHeaderOrigin\enspace}
\renewcommand{\AnswerListHeader}{\textbf{\ExerciseHeaderNB.}\enspace}
\renewcounter{Exercise}[chapter]
\renewcommand{\AtBeginAnswer}{\vspace{-1\baselineskip}}
\newcommand{\pro}{\Exercise}
\newcommand{\qn}{\Question}
\newcommand{\subqn}{\subQuestion}
\newcommand{\sol}{\Answer}

\makeatletter
\renewcommand{\@@@ExeCmd}{%
	\ifnum\@QuestionLevel=0
		\advance \@QuestionLevel by 1
		\begin{list}{\@getExerciseInfo\ExerciseListHeader}%
		{\partopsep \Exepartopsep \labelsep \Exelabelsep \itemsep \Exesep \listparindent 2em%
		\parsep \Exeparsep \topsep \Exetopsep \labelwidth \Exelabelwidth%
		\leftmargin \Exeleftmargin \rightmargin \Exerightmargin}
	\else
		\termineliste{1}\@EndExeBox
	\fi
	\@selectExercise
	\global\@Answerfalse\@BeginExeBox\refstepExecounter%
	\addcontentsline{\ext@exercise}{\toc@exercise}{\ExerciseName\
	\theExercise\ \expandafter{\itshape \ExerciseTitle}\hspace{.66em}}
	\item\ignorespaces\AtBeginExercise
}
\makeatother

\newcounter{soleq}
\counterwithin{soleq}{chapter}
\newenvironment{soleq}{\refstepcounter{soleq}\equation}{\tag{S-\thesoleq}\endequation}

\makeatletter
\newenvironment{problems}{%
	\begin{ExerciseList}
		\global\setbox\temp@Answerbox%
		\vbox{\unvbox\all@Answerbox\vspace{1cm}\section*{参考答案\zhnum{chapter}}\unvbox\@Exercisebox\vskip\z@}%
		\global\setbox\all@Answerbox\copy\temp@Answerbox
		}{
	\end{ExerciseList}
}

\newenvironment{choices}{\begin{compactenum}[\quad\ \,A.{\ \ \ }]}{\end{compactenum}}

\newenvironment{hint}{%
	\ifvmode
		\ignorespaces
	\else
		\quad
	\fi
	\begin{tikzpicture}[baseline=(H.base), every node/.style={signal, draw, very thin, signal to=east, signal from=nowhere, signal pointer angle=120, inner sep=2pt}]
		\node[anchor=mid west] (H) at (0,0) {\heiti\footnotesize 提示};
	\end{tikzpicture}
}{}

\g@addto@macro\frontmatter{\setcounter{tocdepth}{1}
	\setcounter{page}{1}
	\thispagestyle{empty}
	\addtocontents{toc}{\protect\thispagestyle{empty}}
	
	\renewcommand{\contentsname}{目录}
	\tableofcontents
}

\renewcommand{\figurename}{图}
\renewcommand{\tablename}{表}

\renewcommand{\listfigurename}{图片索引}
\renewcommand{\listtablename}{表格索引}

\pretocmd{\listoffigures}{%
	\cleardoublepage
	\phantomsection
	\addcontentsline{toc}{chapter}{\listfigurename}
}{}{}
\pretocmd{\listoftables}{%
	\cleardoublepage
	\phantomsection
	\addcontentsline{toc}{chapter}{\listtablename}
}{}{}

\usepackage[thmmarks, amsmath, hyperref]{ntheorem}
\theoremstyle{change}
\theorembodyfont{\fangsong}
\theoremheaderfont{\sffamily\bfseries\thmheiti}
\theoremseparator{\ }
\setlength{\theorempreskipamount}{2mm}
\setlength{\theorempostskipamount}{2mm}

\newtheorem{theorem}{定理}[chapter]

\theorembodyfont{\songti}
\newtheorem{problem}[theorem]{例题}
\newtheorem{example}[theorem]{例}

\theorembodyfont{}
\theoremstyle{nonumberplain}
\theoremheaderfont{\sffamily\bfseries\CJKfamily{pffont}}
\theoremseparator{\enspace}
\theoremsymbol{}
\newtheorem{solution}{解}

\numberwithin{equation}{section}
\renewcommand{\theequation}{\thesection.\arabic{equation}}
\allowdisplaybreaks

\newcommand{\proy}[1]{\pro[origin={#1}]}

\usepackage{epigraph}

\renewcommand\epigraphflush{flushright}
\renewcommand\epigraphsize{\normalsize}
\setlength\epigraphwidth{0.7\textwidth}

\definecolor{titlepagecolor}{cmyk}{1, 0.6, 0, 0.4}

\newcommand\titlepagedecoration{%
	\begin{tikzpicture}[remember picture,overlay,shorten >= -10pt]
		\coordinate (aux1) at ([yshift=-15pt] current page.north east);
		\coordinate (aux2) at ([yshift=-410pt] current page.north east);
		\coordinate (aux3) at ([xshift=-4.5cm] current page.north east);
		\coordinate (aux4) at ([yshift=-150pt] current page.north east);

		\begin{scope}[titlepagecolor!40, line width=12pt, rounded corners=12pt]
			\draw (aux1) -- coordinate (a) ++(225:5) -- ++(-45:5.1) coordinate (b);
			\draw[shorten <= -10pt] (aux3) -- (a) -- (aux1);
			\draw[opacity=0.6,titlepagecolor,shorten <= -10pt] (b) -- ++(225:2.2) -- ++(-45:2.2);
		\end{scope}
		
		\draw[titlepagecolor, line width=8pt, rounded corners=8pt, shorten <= -10pt] (aux4) -- ++(225:0.8) -- ++(-45:0.8);
		
		\begin{scope}[titlepagecolor!70, line width=6pt, rounded corners=8pt]
			\draw[shorten <= -10pt] (aux2) -- ++(225:3) coordinate[pos=0.45] (c) -- ++(-45:3.1);
			\draw (aux2) -- (c) -- ++(135:2.5) -- ++(45:2.5) -- ++(-45:2.5) coordinate[pos=0.3] (d);   
			\draw (d) -- +(45:1);
		\end{scope}
	\end{tikzpicture}%
}

\begin{document}
	\begin{titlepage}
		\noindent
		{\fontsize{40}{40} \selectfont \bfseries \sffamily \CJKfamily{sectionfont} ICS 问题求解} \par
		\vspace{1cm}
		\epigraph{\LARGE PKU 04832363：计算机系统导论讨论班}{\Large 2021 年秋}
		\null\vfill
		\begin{center} \href{https://www.pku.edu.cn/}{\includegraphics[height=60pt]{logo.eps}} \end{center}
		\titlepagedecoration
	\end{titlepage}
	\newpage
	\thispagestyle{empty}
	\begin{center}
		\LARGE 编译于 {\today} \currenttime
	\end{center}
	\vfill
	\begin{flushleft}
		\large \textbf{汇编：} 王畅 \\
    	\textbf{勘误：} 请写信到 \href{mailto:wchang@pku.edu.cn}{\texttt{wchang@pku.edu.cn}} \\
		\vspace{\baselineskip}
		\textbf{说明：}以下内容是 2021 年秋计算机系统导论讨论班（04832363，班号 16）使用的材料。主要的练习题均来源于往年计算机系统导论课程，编者仅做了汇编和微小的改动。考题只选择了 2018 年及之前的内容，不选用更新的题目是因为 2019 年和 2020 年的试卷将留给大家复习备考时整套使用。有关 Lab 的指导则为编者自行组织。后面附有部分题目的分析和答案，“部分”选择的要旨是对较难的题目作详细解释；对于往年题我们只给出较难的题或者无答案题的分析。
	\end{flushleft}

	\frontmatter

	\mainmatter
	\chapter{位级表示}
	\begin{summary}
		\begin{compactitem}
			\item 知道整数、浮点数在系统的存储方式（字节序）。
			\item 熟练掌握整数、浮点数的位级表示规则，快速完成其同十进制数的相互转换。
			\item 理解整数、浮点数规范中的各类特殊数的计算方法及其性质。
			\item 运用浮点数舍入的规则进行运算，知道类型转换的基本规则，能针对整数、浮点数的一些常见“反常情况”进行判断。
		\end{compactitem}
	\end{summary}

	\begin{problems}
		\pro 在 x86-64 机器上，定义 \texttt{unsigned int A = 0x123456}。请画出 \texttt{A} 在内存中的存储方式：
		\begin{table}[H]
			\centering
			\begin{tabular}{|c|c|c|c|c|c|c|c|}
				\hline
				... & 低地址 & \multicolumn{4}{c|}{\texttt{A}} & 高地址 & ... \\ \hline
				\multicolumn{2}{|c|}{...} & {\qquad \qquad} & {\qquad \qquad} & {\qquad \qquad} & {\qquad \qquad} & \multicolumn{2}{c|}{...} \\ \hline
			\end{tabular}
		\end{table}
		定义 \texttt{unsigned short B[2] = \{0x1234, 0x5678\}}。请画出 \texttt{B} 在内存中的存储方式：
		\begin{table}[H]
			\centering
			\begin{tabular}{|c|c|c|c|c|c|c|c|}
				\hline
				... & 低地址 & \multicolumn{4}{c|}{\texttt{B}} & 高地址 & ... \\ \hline
				\multicolumn{2}{|c|}{...} & {\qquad \qquad} & {\qquad \qquad} & {\qquad \qquad} & {\qquad \qquad} & \multicolumn{2}{c|}{...} \\ \hline
			\end{tabular}
		\end{table}
		\pro 在 x86-64 机器上，有下列 C 代码：
		\begin{minted}[frame=single, fontsize=\small]{c}
    int main() {
        unsigned int A = 0x11112222;
        unsigned int B = 0x33336666;
        void *x = (void *)&A;
        void *y = 2 + (void *)&B;
        unsigned short P = *(unsigned short *)x;
        unsigned short Q = *(unsigned short *)y;
        printf("0x%04x", P + Q);
        return 0;
    }
		\end{minted}
		运行该代码，结果是什么？
		\pro 在 x86-64 机器上，有下列 C 代码：
		\begin{minted}[frame=single, fontsize=\small]{c}
    int main() {
        char A[12] = "11224455";
        char B[12] = "11445577";
        void *x = (void *)&A;
        void *y = 2 + (void *)&B;
        unsigned short P = *(unsigned short *)x;
        unsigned short Q = *(unsigned short *)y;
        printf("0x%04x", Q - P);
        return 0;
    }
		\end{minted}
		运行该代码，结果是什么？
		\pro 在 x86-64 机器上，有如下的定义：
		\begin{minted}[frame=single, fontsize=\small]{c}
    int x = ...; // 表达式 A
    int y = ...; // 表达式 B
    unsigned int ux = x;
    unsigned int uy = y;
		\end{minted}
		判断下表中的表达式是否等价：
		\begin{table}[H]
			\centering
			\begin{tabular}{|c|c|c|}
				\hline
				序号 & 表达式 \verb|A| & 表达式 \verb|B| \\ \hline
				1 & \verb|x > y| & \verb|ux > uy| \\ \hline
				2 & \verb+(x > 0) || (x < ux)+ & \verb|1| \\ \hline
				3 & \verb|x ^ y ^ x ^ y ^ x| & \verb|x| \\ \hline
				4 & \verb|((x >> 1) << 1) <= x| & \verb|1| \\ \hline
				5 & \verb|((x / 2) * 2) <= x| & \verb|1| \\ \hline
				6 & \verb|x ^ y ^ (~x) - y| & \verb|y ^ x ^ (~y) - x| \\ \hline
				7 & \verb|(x == 1) && (ux - 2 < 2)| & \verb|(x == 1) && ((!!ux) - 2 < 2)| \\ \hline
			\end{tabular}
		\end{table}
		\begin{hint}
			减法的运算优先级比按位异或高。布尔运算的结果都是有符号数。
		\end{hint}
		\pro 下列代码的目的是将字符串 \texttt{A} 的内容复制到字符串 \texttt{B}，覆盖 \texttt{B} 原有的内容，并输出“Hello World”；但实际运行输出是“Buggy Codes”。尝试找到代码中的错误。
		\begin{minted}[frame=single, fontsize=\small]{c}
    int main() {
        char A[12] = "Hello World";
        char B[12] = "Buggy Codes";
        int pos;
        for (pos = 0; pos - sizeof(B) < 0; pos++)
            B[pos] = A[pos];
        printf("%s\n", B);
    }
		\end{minted}
		\pro 假设某浮点数格式为 1 位符号、3 位阶码、4 位小数。下表给出了用该格式表达的浮点数 $(-1)^SM \cdot 2^E$ 与其二进制表示的关系。完成下表。
		\begin{table}[H]
			\centering
			\begin{tabular}{|c|c|c|c|c|}
				\hline
				描述 & 二进制表示 & $M$（写成分数） & $E$ & $f$ \\ \hline
				负零 &  & / & / & \texttt{-0.0} \\ \hline
				/ & \verb|01000101| &  &  &  \\ \hline
				最小的非规格化负数 &  &  &  &  \\ \hline
				最大的规格化正数 &  &  &  &  \\ \hline
				一 &  &  &  & \texttt{1.0} \\ \hline
				/ &  &  &  & \texttt{5.5} \\ \hline
				$+\infty$ &  & / & / & / \\ \hline
			\end{tabular}
		\end{table}
		\pro 假设浮点数格式 A 为 1 位符号、3 位阶码、4 位小数，浮点数格式 B 为 1 位符号、4 位阶码、3 位小数。回答下列问题。
			\qn 格式 A 中有多少个二进制表示对应于正无穷大？
			\qn 考虑能精确表示的实数的最大绝对值。A 比 B 大还是比 B 小，还是两者一样？
			\qn 考虑能精确表示的实数的最小非零绝对值。A 比 B 大还是比 B 小，还是两者一样？
			\qn 考虑能精确表示的实数的个数。A 比 B 多还是比 B 少，还是两者一样？
		\pro 判断下列说法的正确性。
			\qn 对于任意的单精度浮点数 \texttt{a} 和 \texttt{b}，如果 \texttt{a > b}，那么 \texttt{a + 1 > b}。
			\qn 对于任意的单精度浮点数 \texttt{a} 和 \texttt{b}，如果 \texttt{a > b}，那么 \texttt{a + b > b + b}。
			\qn 对于任意的单精度浮点数 \texttt{a} 和 \texttt{b}，如果 \texttt{a > b}，那么 \texttt{a + 1 > b + 1}。
			\qn 对于任意的双精度浮点数 \texttt{d}，如果 \texttt{d < 0}，那么 \texttt{d * d > 0}。
			\qn 对于任意的双精度浮点数 \texttt{d}，如果 \texttt{d < 0}，那么 \texttt{d * 2 < 0}。
			\qn 对于任意的双精度浮点数 \texttt{d}，\texttt{d == d}。
			\qn 将 \texttt{float} 转换成 \texttt{int} 时，既有可能造成舍入，又有可能造成溢出。
		\pro 遵循 IEEE 754 浮点数标准，考虑下列代码：
		\begin{minted}[frame=single, fontsize=\small]{c}
    for (int x = 0; ; x++) {
        float f = x;
        if (x != (int)f) {
            printf("%d", x);
            break;
        }
    }
		\end{minted}
		试问代码的运行结果是什么？或者死循环？
		\pro 遵循 IEEE 754 浮点数标准，考虑下列代码：
		\begin{minted}[frame=single, fontsize=\small]{c}
    int x = 33554466; // 2^25 + 34
    int y = x + 8;
    for ( ; x < y; x++) {
        float f = x;
        printf("%d ", x - (int)f);
    }
		\end{minted}
		写出程序的运行结果。
	\end{problems}

\chapter{位级表示{---}往年考题}
	\begin{problems}
		\proy{2018} 下列哪种类型转换既可能导致溢出，又可能导致舍入？
		\begin{choices}
			\item \texttt{int} 转 \texttt{float}
			\item \texttt{float} 转 \texttt{int}
			\item \texttt{int} 转 \texttt{double}
			\item \texttt{float} 转 \texttt{double}
		\end{choices}
		\proy{2018} 在采用小端法存储机器上运行下面的代码，输出的结果是？（\texttt{int}、\texttt{unsigned} 为 32 位长，\texttt{short} 为 16 位长，0\textasciitilde9 的 ASCII 码是 \texttt{0x30}\textasciitilde\texttt{0x39}）
		\begin{minted}[frame=single, fontsize=\small]{c}
    char *s = "2018";
    int *p1 = (int *)s;
    short s1 = (*p1) >> 12;
    unsigned u1 = (unsigned) s1;
    printf("0x%x\n", u1);
		\end{minted}
		\begin{choices}
			\item \texttt{0x00002303}
			\item \texttt{0x00032303}
			\item \texttt{0xffff8313}
			\item \texttt{0x00008313}
		\end{choices}
		\proy{2018} 考虑如下函数
		\begin{minted}[frame=single, fontsize=\small]{c}
    void XOR(int x, int y) {
        y = x ^ y;
        x = x ^ y;
        y = x ^ y;
        printf(x, y);
    }
		\end{minted}
		则 $\texttt{XOR}(a, b)$ 的输出结果是什么？
		\begin{choices}
			\item $a, b$
			\item $b, a$
			\item $b, 0$
			\item $b, a\ \verb|^|\ b$
		\end{choices}
		\proy{2017} 假定一个特殊设计的计算机，将 \texttt{int} 型数据的长度从 4 字节 扩展为 $4N$ 字节，采用大端法。现将该 \texttt{int} 型所能表示的最小负数写入内存中，如下图所示。其中每个小矩形代表一个字节，请问 X 位置这个字节中的值是多少？
		\begin{figure}[H]
			\centering
			\MemoryLayout{%
				8/white/\relax,
				24/gray!25/{$4N$ 字节},
				32/white/\relax
			}{%
				\draw[thick] (0,0) -- ++(0,3) node[above] {低地址};
				\draw[thick] (-32,0) -- ++(0,3) node[above] {高地址};
				\draw[-{Stealth[angle'=45]}, thick] (-8.5,1) -- ++(0,3) node[above] {X};
			}
		\end{figure}
		\begin{choices}
			\item \texttt{00000000}
			\item \texttt{01111111}
			\item \texttt{10000000}
			\item \texttt{11111111}
		\end{choices}
		\proy{2017} 以下说法正确的是：
		\begin{choices}
			\item 负数加上负数结果都为负数
			\item 正数加上正数结果都为正数
			\item 用 \verb|&| 和 \verb|~| 可以表示所有的逻辑与或非操作
			\item 用 \verb|&| 和 \verb+|+ 可以表示所有的逻辑与或非操作
		\end{choices}
		\proy{2017, 2016} 若我们采用基于 IEEE 浮点格式的浮点数表示方法，阶码字段（exp）占据 $k$ 位，小数字段（frac）占据 $n$ 位，则最小的规格化正数是：
		\begin{choices}
			\item $(1-2^{-n}) \cdot 2^{-2^{k-1}+2}$
			\item $2^{-2^{k-1}+2}$
			\item $2^{-n} \cdot 2^{-2^{k-1}+2}$
			\item $(1-2^{-n}) \cdot 2^{-2^k+1}$
		\end{choices}
		\proy{2016} 假定编译器规定 \texttt{int} 和 \texttt{short} 型长度分别为 32 位和 16 位，执行下列语句：\texttt{unsigned short x = 65530; unsigned int y = x;}，得到 \texttt{y} 的机器数是 \rule{2.5cm}{0.25mm}。（用 16 进制表示，勿省略前导的 0）
		\proy{2016} 一个 C 语言程序在一台 32 位机器上运行。程序中定义了三个变量 $x, y, z$，其中 $x$ 和 $z$ 为 \texttt{int} 型，$y$ 为 \texttt{short} 型。当 $x=127, y=-9$ 时，执行赋值语句 $z=x+y$ 后，$z$ 的值是 \rule{2.5cm}{0.25mm}。（用 16 进制表示，勿省略前导的 0）
		\proy{2016} 若按 IEEE 浮点标准的单精度浮点数（符号位 1 位，阶码字段 exp 占据 8 位，小数字段 frac 占据 23 位）表示 $-8.25$，结果是 \rule{2.5cm}{0.25mm}。（用 16 进制表示）
		\proy{2015} 给定一个实数，会因为该实数表示成单精度浮点数而发生误差。不考虑 \texttt{NaN} 和 \texttt{Inf} 的情况，该绝对误差的最大值为：
		\begin{choices}
			\item $2^{103}$
			\item $2^{104}$
			\item $2^{230}$
			\item $2^{231}$
		\end{choices}
		\proy{2016} 现有一个二进制浮点的表示规则，其中 $E$ 为指数部分（3 比特），bias为 3；$M$ 为小数部分（5 比特），采用二进制补码表示形式，且取值 $0.5 \leq |M|<1$，$s$ 是浮点的符号位。该形式包含一个值为 1 的隐藏位。问 $+5_{10}$ 在该表示下的值是下列哪一个？
		\begin{choices}
			\item \texttt{010001100}
			\item \texttt{010100100}
			\item \texttt{011011010}
			\item \texttt{011110101}
		\end{choices}
		\proy{2015} 在 32 位平台上，按 C90 标准判定以下语句中结果为假的是：
		\begin{choices}
			\item \verb|return INT_MIN < INT_MAX;|
			\item \verb|return -2147483648 < 2147483647;|
			\item \verb|int a = -2147483648; return a < 2147483647;|
			\item \verb|return -2147483647 - 1 < 2147483647;|
		\end{choices}
		\begin{hint}
			C90 中的类型转换顺序：\texttt{int -> long -> unsigned -> unsigned long}，$2^{31}=2147483648$。
		\end{hint}
		\proy{2015} 关于浮点数，以下说法正确的是：
		\begin{choices}
			\item 给定任意浮点数 $a, b, x$，如果 $a>b$ 成立(指求值为 1)，则一定有 $a+x>b+x$ 成立
			\item 给定任意浮点数 $a, b, x$，如果 $a>b$ 不成立(指求值为 0)，则一定有 $a+x>b+x$ 不成立
			\item 不考虑结果为 \texttt{NaN, Inf} 或运算过程发生溢出的情况，高精度浮点数一定得到比低精度浮点数更精确或相同的结果
			\item 不考虑结果为 \texttt{NaN, Inf} 的情况，高精度浮点数一定得到比低精度浮点数更精确或相同的结果
		\end{choices}
		\proy{2014} 设整数均为 32 位，其中 \texttt{unsigned x = 0x00000001; int y = 0x80000000; int z = 0x80000001;}。以下表达式求值为 1 的是：
		\begin{choices}
			\item \verb|(-1) < x|
			\item \verb|(-y) > -1|
			\item \verb|~y + y == -1|
			\item \verb|(z << 4) > (z * 16)|
		\end{choices}
		\proy{2014} 下面说法正确的是：
		\begin{choices}
			\item 数 0 的反码表示是唯一的
			\item 数 0 的补码表示不是唯一的
			\item \verb|1000_1111_1110_1111_1100_0000_0000_0000| 表示的唯一整数是 \texttt{0x8FEFC000}
			\item \verb|1000_1111_1110_1111_1100_0000_0000_0000| 如果是单精度浮点数，则其值是 $-(1.110111111)_2 \cdot 2^{31-127}$
		\end{choices}
		\proy{2016, 2014} 下面表达式中为真的是：
		\begin{choices}
			\item \verb|(unsigned)-1 < -2|
			\item \verb|2147483647 > (int)2147483648U|
			\item \verb|(0x80005942 >> 4) == 0x09005942|
			\item \verb|2147483647 + 1 != 2147483648|
		\end{choices}
        \proy{2014} 下面关于 IEEE 浮点数标准说法正确的是哪个？
		\begin{choices}
			\item 在位数一定的情况下，不论怎么分配 exponent bits 和 fraction bits，所能表示的数的个数是不变的。
			\item 若甲类浮点数有 10 位，乙类浮点数有 11 位，那么甲所能表示的最大数一定比乙小。
			\item 若甲类浮点数有 10 位，乙类浮点数有 11 位，那么甲所能表示的最小正数一定比乙小。
		    \item “\texttt{0111000}”可能是 7 位浮点数的 \texttt{NaN} 表示。
		\end{choices}
		\proy{2014} 假设有下面 $x$ 和 $y$ 的程序定义：
		\begin{minted}[frame=single, fontsize=\small]{c}
    int x = a >> 2;
    int y = (x + a) / 4;
		\end{minted}
		那么有 \rule{2.5cm}{0.25mm} 个位于闭区间 $[-8, 8]$ 内的整数 $a$ 能使得 $x$ 和 $y$ 相等。
        \proy{2013} 对于 IEEE 浮点数，如果减少 1 位指数位，将其用于小数部分，下列叙述正确的是哪个？
		\begin{choices}
			\item 能表示更多数量的实数值，但实数值取值范围比原来小了。
			\item 能表示的实数数量没有变化，但数值的精度更高了。
			\item 能表示的最大实数变小，最小的实数变大，但数值的精度更高。
			\item 以上说法都不正确。
        \end{choices}
		\proy{2017} 考虑有一种基于 IEEE 浮点格式的 9 位浮点表示格式 A。格式 A 有 1 个符号位、$k$ 个阶码位、$n$ 个小数位。现在已知 $-9$ 的位模式可以表示为 \texttt{101100010}。回答以下问题。（注：阶码偏移量为 $2^{k-1}-1$）
			\qn 求 $k$ 和 $n$ 的值。
			\qn 基于格式 A，请填写下表。值的表示可以写成整数（如 16），或者写成分数（如 17/64）。
			\begin{table}[H]
				\centering
				\begin{tabular}{|c|c|c|}
					\hline
					描述 & 二进制表示 & 值 \\ \hline
					最大的非规格化数 & {\qquad \qquad} & {\qquad \qquad} \\ \hline
					最小的正规格化数 & {\qquad \qquad} & {\qquad \qquad} \\ \hline
					最大的规格化数 & {\qquad \qquad} & {\qquad \qquad} \\ \hline
				\end{tabular}
			\end{table}
			\qn 假设格式 A 变为 1 个符号位、$k+1$ 个阶码位、$n-1$ 个小数位，那么能表示的实数数量会怎样变化？数值的精度会怎样变化？（回答增加、降低或不变即可）
		\proy{2016} 在 64 位机器上，判断表达式是否对代码生成的变量恒成立（指求值为 1）。
		\begin{minted}[frame=single, fontsize=\small]{c}
    /* random_int() 函数返回一个随机的 int 类型值 */
    int x = random_int();
    int y = random_int();
    int z = random_int();
    unsigned ux = (unsigned)x;
    long lx = (long)x; /* long 为 64 位 */
    long ly = (long)y;
    double dx = (double)x;
    double dy = (double)y;
    double dz = (double)z;
		\end{minted}
			\qn \verb+(x >= 0) || (3 * x < 0)+
			\qn \verb+(x >= 0) || (x < ux)+
			\qn \verb|((x >> 1) << 1) <= x|
			\qn \verb|((x - y) << 3) + (x >> 1) - y == 8 * x - 9 * y + x / 2|
			\qn \verb|(x - y > 0) == ((y + ~x + 1) >> 31 == 1)|
			\qn \verb|dx + dy == (double) (y + x)|
			\qn \verb|dx + dy + dz == dz + dy + dx|
			\qn \verb|(int)((lx + ly) >> 1) == ((x & y) + ((x ^ y) >> 1))|
		\proy{2016} 假设 C 语言中新定义了一种数据类型 T, 该类型为 12-bit 长的浮点数，此浮点数遵循 IEEE 浮点数格式，其字段划分如下：符号位(s) 1-bit、阶码字段(exp) 6-bit、小数字段(frac) 5-bit。
			\qn 若将该格式下能表示的所有正规格数从小到大依次排列，则相邻两数之间差值的最小值为 \rule{2.5cm}{0.25mm}，最大值为 \rule{2.5cm}{0.25mm}。
			\qn 现定义了如下变量：
			\begin{minted}[frame=single, fontsize=\small]{c}
    T a = -15.875;
    T b = (1 << 28) + (1 << 24) + (1 << 22);
    T c = a * b;
			\end{minted}
			给出各变量的二进制表示。
		\proy{2015} 对于下面的每一个表达式，请选择以下选项中的\emph{一个或多个}，使得该表达式恒成立，如果没有满足条件的选项则填写 none：A. \texttt{<} \quad B. \texttt{>} \quad C. \texttt{==} \quad D. \texttt{!=} \quad E. none。
		
		题目中出现的变量定义如下（浮点数保证不是 \texttt{NaN} 或者 \texttt{Inf}）：\texttt{int x, y; unsigned ux = x; double d;}。
		\qn 如果 \texttt{x > 0}，则 \texttt{x + 1} \rule{2.5cm}{0.25mm} \texttt{0}
		\qn 如果 \texttt{x > y}，则 \texttt{ux} \rule{2.5cm}{0.25mm} \texttt{y}
		\qn 如果 \verb|((x << 31) >> 31) < 0|，则 \verb|x & 1| \rule{2.5cm}{0.25mm} \texttt{0}
		\qn 如果 \verb|((unsigned char)x >> 1) < 64|，则 \texttt{(char)x} \rule{2.5cm}{0.25mm} \texttt{0}
		\qn 如果 \texttt{d < 0}，则 \texttt{d * 2} \rule{2.5cm}{0.25mm} \texttt{0}
		\qn 如果 \texttt{d < 0}，则 \texttt{d * d} \rule{2.5cm}{0.25mm} \texttt{0}
		\qn \verb|x ^ y ^ (~x) - y| \rule{2.5cm}{0.25mm} \verb|y ^ x ^ (~y) - x|
		\qn \verb|(((!!ux)) << 31) >> 31)| \rule{2.5cm}{0.25mm} \verb|(((!!x) << 31) >> 31)|
		\proy{2015} 考虑一种 12-bit 长的浮点数，此浮点数遵循 IEEE 浮点数格式，浮点数的字段划分如下：符号位(s) 1-bit、阶码字段(exp) 4-bit、小数字段(frac) 7-bit。回答下列问题。
			\qn 请写出在下列区间中包含多少个用上面规则精确表示的浮点数：$[1, 2); [2, 3)$。
			\qn 填写下表。
			\begin{table}[H]
				\centering
				\begin{tabular}{|c|c|}
					\hline
					描述 & 二进制表示 \\ \hline
					最大的非规格化数 & {\qquad \qquad \qquad \qquad} \\ \hline
					最小的正规格化数 & {\qquad \qquad \qquad \qquad} \\ \hline
					$17 \frac{1}{16}$ & {\qquad \qquad \qquad \qquad} \\ \hline
					$-\frac{1}{8192}$ & {\qquad \qquad \qquad \qquad} \\ \hline
					$20 \frac{3}{8}$ & {\qquad \qquad \qquad \qquad} \\ \hline
					$-\infty$ & {\qquad \qquad \qquad \qquad} \\ \hline
				\end{tabular}
			\end{table}
		\proy{2014} 回答下列问题。
			\qn 假设下列 \texttt{unsigned} 和 \texttt{int} 数均为 5 位（有符号整型用补码运算表示）：\texttt{int y = - 7; unsigned z = y;}。确定 \texttt{y, z} 和最小有符号整数的十进制表示和二进制表示。
			\qn 按照 IEEE 浮点数标准，首先将下列两个数表示 $(-1)^sM \cdot 2^E$ 的形式，然后写出其二进制表示：$0.375, -12.5$。
	\end{problems}
	\chapter{Datalab 热身}
    \begin{help}
        \begin{compactitem}
            \item 在 ANSI C 中，变量的声明必须出现在每个花括号包容的 scope 的开头，否则会发生编译错误。
            \item 在使用 \texttt{int, long} 的转换时需要小心，例如立即数默认为 \texttt{int}，经 \texttt{!} 操作后的整数类型都会变为 \texttt{int}，等等。
            \item 可能会有些方法刻意运用整型溢出的操作（以节省操作数），正常完成一般用不到。如果你要使用，请特别注意整型溢出中的编译器行为。
            \item 通常来说，datalab 的完成不会遇到太多困难。如果你在整数部分遇到了一些麻烦，可以看一看下面的例题。这里的例子给出的提示都比较充分，展现了一些处理技巧。这些例子做不出可以找同学或者我索取答案。
            \item 浮点部分的题目允许使用的操作大大增多，所以用暴力方法就能完成，主要考察的是大家对 IEEE 754 标准的理解。（其实大多数整数部分的题目也是这样。）
            \item 有两本有趣的参考书：\textit{Hacker's Delight} 和 \textit{Matters Computational: Ideas, Algorithms, Source Code}，以及一个链接 \url{https://graphics.stanford.edu/~seander/bithacks.html}，里面介绍了一些位运算魔法。Datalab 和往年有一些比较过分的考题大量出自前一本书。但是，做完 Lab 之前\uwave{绝对不要}去看！！！否则会被视为抄袭。 
        \end{compactitem}
    \end{help}

    总地来说，做 datalab 时，不要总是惦记着只能用某些运算，而是先要宏观地找出方法和思路。对于中等难度的题，一般都可以用“拼凑”的方法，先算出什么...再算出什么... 使用模块化的思考方式，例如一些常见的操作可以总结出来随时使用，等等。
    
    即便是遇到很难的题，直接用（不加限制的） C 代码写出算法都是很容易的，你可以先这样做，然后设法改写成符合要求的代码。如有必要，你可以直接展开循环和条件表达式。

    我们下面看几个例子。本学期的 Lab 内容大多数和它们不同，而且这里都没有给出答案，因此不要期望在其中找到答案而不思考。但这里基本把需要用到的思路都涵盖了，因而你如果不想因为被变相剧透思路而失去做题的乐趣，我强烈建议你现在就关掉本文档。

    \begin{example}[de Morgan 律]
        我们知道 \verb+~, &, |+ 可以表达“所有”的逻辑表达式，比方说异或 \verb+x ^ y = (x & ~y) | (~x ^ y)+。而实际上 \verb+~, |+ 就足以做到这件事，因为我们有 \verb+x & y = ~(~x | ~y)+。这样的情况称为\emph{连接词的完备集}，以后大家还会反复学到。

        当然，其实单纯一个与非或者一个或非也够了。假设\emph{或非} 用 $\downarrow$ 表示（即 $\texttt{x} \downarrow \texttt{y} = \verb+~(x | y)+$），请你完成以下代码：唯一允许的位运算操作符是 $\downarrow$。
        \begin{minted}[frame=single, fontsize=\small]{c}
    unsigned xor_with_nor(unsigned x, unsigned y) {
        return _________________________; // return x ^ y with only NOR
    }
        \end{minted}
    \end{example}

    \begin{example}[表达式]
        试用位级运算的技术计算表达式 \texttt{cond ? t : f}。根据注释中的提示尝试补全下面的代码，假设 \texttt{cond} 的输入总是 1 或者 0。
        \begin{minted}[frame=single, fontsize=\small]{c}
    int conditional(int cond, int t, int f) {
        /* Compute a mask that equals 0x00000000
         * or 0xFFFFFFFF depending on the value of cond */
        int mask = ____________________;
        /* Use the mask to toggle between returning t or returning f */
        return ____________________;
    }
        \end{minted}
        这个例子的用处是，如果你一定想要用 \texttt{if}，可以用这个办法避开 datalab 的限制。更显然地，减法、常数乘法都是事实上可以用的，此外你也可以思考如何表达 \texttt{==}。
        \begin{minted}[frame=single, fontsize=\small]{c}
    int equal(int x, int y) {
        return _________________________; // return x == y
    }

    int minus(int x, int y) {
        return _________________________; // return x - y
    }
        \end{minted}
    \end{example}

    \begin{example}[模拟操作 1]
        你可能以前听说过，如果要计算两个无符号数的平均 $\left \lfloor \frac{x+y}{2} \right \rfloor$ 而不发生溢出可以用：\verb|(x & y) + ((x ^ y) >> 1)|。不难看出其原理：\verb|x ^ y| 是半加（不考虑进位）的部分，不会发生溢出，右移即可；另一部分则包括进位，当且仅当两位均为 1 时发生，而如果将这个进位保留在原地，恰好就是除以 2 的效果。这样的思想来自于计算机的加法器。
        \begin{figure}[H]
            \centering
            \tikzstyle{branch}=[fill, shape=circle, minimum size=3pt, inner sep=0pt]
            \begin{tikzpicture}[label distance=2mm]
                \node (x) at (-1,6) {$x$};
                \node (y) at ($(x) + (0,-1.2)$) {$y$};
                \node[not gate US, draw] at ($(x)+(0.5,-0.8)$) (notx) {};
                \node[not gate US, draw] at ($(y)+(0.5,-0.8)$) (noty) {};
                \node[and gate US, draw, rotate=-90, logic gate inputs=nn] at (1,3) (A) {};
                \node[and gate US, draw, rotate=-90, logic gate inputs=nn] at ($(A)+(2,0)$) (B) {};
                \node[and gate US, draw, rotate=-90, logic gate inputs=nn] at ($(B)+(2,0)$) (C) {};
                \node[or gate US, draw, rotate=-90, logic gate inputs=nn] at ($(A)+(1,-1.5)$) (D) {};
                \foreach \i in {x,y} {
                    \path (\i) -- coordinate (punt\i) (\i |- not\i.input);
                    \draw (\i) |- (punt\i) node[branch] {} |- (not\i.input);
                }
                \draw (puntx) -| (C.input 1);
                \draw (punty) -| (C.input 2);
                \draw (puntx) -| (B.input 1);
                \draw (punty) -| (A.input 2);
                \draw (notx) -| (A.input 1);
                \draw (noty) -| (B.input 2);
                \draw (A.output) -- ([yshift=-0.2cm] A.output) -| (D.input 2);
                \draw (B.output) -- ([yshift=-0.2cm] B.output) -| (D.input 1);
                \draw (C) -- ($(C) + (0, -1.8)$) -- node[right] {$C$} ($(C) + (0, -2.5)$);
                \draw (D.output) -- node[right] {$S$} ($(D) + (0, -1)$);
            \end{tikzpicture}
            \caption{半加器的结构，它计算两个一位整数的和 $S$ 及其进位 $C$。可以看出 $S = x \oplus y, C=x \& y$。}
        \end{figure}

        一般来说，有符号数的处理比无符号数稍微需要一些讨论。根据上面的方法，请你尝试给出两个有符号数的平均 $\left \lfloor \frac{x+y}{2} \right \rfloor$ 而不发生溢出的算法；对 $\left \lceil \frac{x+y}{2} \right \rceil$ 做同样的事。
        
        与此相关，请问如何用位运算检查 $x+y$ 是否溢出了？\begin{hint} 检查符号位。 \end{hint}
        \begin{minted}[frame=single, fontsize=\small]{c}
    int addOK(int x, int y) {
        ________________________________;
        ________________________________;
        // You can add more lines
        return ____; // Determine if we can compute x + y without overflow
    }
        \end{minted}
    \end{example}

    \begin{example}[计数]
        让我们来决定一个给定数的二进制表示中，1 的个数是偶数还是奇数；奇数返回 1，否则返回 0。允许使用 datalab 整数部分规定的所有运算符。
        
        这里我们采用\emph{分治}的策略，首先考察比较短的数。例如，当所处理的数只有两位时，采用的代码十分显然：
        \begin{minted}[frame=single, fontsize=\small]{c}
    int bitParity2bit(int x) {
        int bit1 = 0b01 & x;
        int bit2 = 0b01 & (x >> 1);
        return bit1 ^ bit2;
    }
        \end{minted}
        将异或理解为 $\mathbb F_2$ 上的加法是非常有益的。

        对于四位整数，我们两位两位操作，并让操作的过程某种意义上“并行”进行。首先分别确定 1、2 和 3、4 位的奇偶性，所得结果再设法异或一次。（思考：\texttt{mask2} 可以改成其他数吗？）
        \begin{minted}[frame=single, fontsize=\small]{c}
    int bitParity4bit(int x) {
        int mask = 0b0101;
        int halfParity = (mask & x) ^ (mask & (x >> 1));
        int mask2 = 0b0011;
        return (mask2 & halfParity) ^ (mask2 & (halfParity >> 2));
    }
        \end{minted}

        现在，请你尝试根据上面的提示，补全下面针对八位整数的算法（要求使用操作符数目不超过 12 个）。
        \begin{minted}[frame=single, fontsize=\small]{c}
    int bitParity8bit(int x) {
        int mask = ____________________;
        int quarterParity = ____________________;
        int mask2 = ____________________;
        int halfParity = ____________________;
        int mask3 = ____________________;
        return ____________________;
    }
        \end{minted}
        最后，试完成针对十六位或更长整数的算法（注意，根据 datalab 要求，立即数有效长度不能超过 8 位）。
        
        作为模仿练习，请你对问题“决定一个给定数的二进制表示中有多少个 1”做同样的事。
        \begin{minted}[frame=single, fontsize=\small]{c}
    /* Let's count how many bits are set in a number */
    int bitCount8bit(int x) {
        int mask = ____________________;
        int quarterSum = ____________________;
        int mask2 = ____________________;
        int halfSum = ____________________;
        int mask3 = ____________________;
        return ____________________ + ____________________;
    }
        \end{minted}
    \end{example}

    \begin{example}[模拟操作 2]
        对于一个无符号整数，我们希望将它的位反过来（记为 $\operatorname{rev}(\cdot)$），例如 \texttt{0x01234567} 倒过来就是 \texttt{0xE6A2C480}。

        这个某种意义上是一个广为流传的\sout{面试题}。我们沿用\emph{分治}的思想。假设二进制数 $x$ 有 $2m$ 位 $(\underbracket{b_{2m-1} \dotsm b_m}_{x_h} \underbracket{b_{m-1} \dotsm b_0}_{x_l})_2$，那么显然有
        \[ \operatorname{rev}(x) = \operatorname{rev}(x_l) \operatorname{rev}(x_h). \]
        根据这个思路，请你补全以下代码：
        \begin{minted}[frame=single, fontsize=\small]{c}
    unsigned reverseBits(unsigned n) {
        n = (n >> 16) | (n << 16);
        n = ((n & __________) >> 8) | ((n & __________) << 8);
        n = ________________________ | _______________________;
        n = __________________________________________________;
        n = __________________________________________________;
        return n;
    }
        \end{minted}

        读过两个分治的操作之后，请思考：给一个整数，如何计算它有多少个前导零？允许的操作符包括所有的位级运算，以及加法和 \texttt{!}。\begin{hint} 二分查找。 \end{hint}
    \end{example}

    \begin{example}[算术运算]
        考虑这样一个问题：对一个无符号整数 $x$, 计算 $x \bmod 5$。除了加法之外不能用其他算术操作。

        方法上这很标准，假设 $x=(b_{31} \dotsm b_1b_0)_2$，通过计算 $2^i \pmod 5$ 的周期，我们知道
        \[ x \equiv \sum_{i=0}^{31} b_i2^i \equiv 1b_0+2b_1+4b_2+3b_3+1b_4 + \dotsb + 3b_{31} \pmod 5. \]
        技术上，对一个四位无符号整数 $(b_3b_2b_1b_0)_2$ 计算 $b_0+2b_1+4b_2+3b_3$ 可以非常直接：比如首先取得各位，然后用左移配合加法计算结果。你也可以思考有什么更省操作符数目的方法。
    \end{example}
	\chapter{汇编语言}
    \begin{summary}
        \begin{compactitem}
            \item 知道 x86-64 中指令、程序计数器、寄存器、条件码、内存等概念，记住重要的寄存器名称、含义和使用规范，清楚掌握数据传送指令和操作数的正确用法。
            \item 掌握条件分支、条件传送、各种循环、跳转表的翻译方式，能熟练地在汇编语言中识别控制流结构，快速完成机器码、汇编语言、C 代码的相互转换。
            \item 理解 x86-64 系统栈空间的分布和管理方式，知道过程调用中的重要寄存器和相关保存指令，会复述过程调用的整个过程并绘制栈空间的变化情况。
            \item 知道结构体、联合体在内存中的存储情况，掌握用对齐规则访问结构体和联合体，以及计算其实际大小的方法。会处理复杂的指针和函数指针问题。
        \end{compactitem}
    \end{summary}

    \begin{problems}
        \pro 判断下列 x86-64 ATT 操作数格式是否合法。
            \qn \verb|8(%rax, , 2)|
            \qn \verb|$30(%rax, %rax, 2)|
            \qn \verb|0x30|
            \qn \verb|13(, %rdi, 4)|
            \qn \verb|(%rsi, %rdi, 6)|
            \qn \verb|%ecx|
            \qn \verb|(%ecx)|
            \qn \verb|(%rbp, %rsp)|
        \sol 1 不正确；2 不正确，\verb|$| 只用来表示立即数；3 正确，是内存地址 \verb|0x30|；4 正确；5 不正确，缩放比例只能是 1、2、4、8；6 正确；7 不正确，x86-64 不允许将除了 64 位寄存器以外的寄存器作为寻址模式基地址；8不正确，\verb|%rsp|不能作为操作数！
        \pro 假设 \verb|%rax|、\verb|%rbx| 的初始值都是 0。根据下列一段汇编代码，写出每执行一步后两个寄存器的值。
        \begin{minted}[frame=single, fontsize=\small]{gas}
    movabsq   $0x0123456789ABCDEF, %rax
    movw      %ax, %bx
    movswq    %bx, %rbx
    movl      %ebx, %eax
    movabsq   $0x0123456789ABCDEF, %rax
    cltq
        \end{minted}
        \sol 注意 \verb|movl| 会清零高 32 位。如果 \verb|movabsq| 的立即数是负数，主要要取绝对值（本题中不需要）。答案为：
        \begin{compactenum}
            \item \ \verb|0x0123456789ABCDEF, 0x0000000000000000|
            \item \ \verb|0x0123456789ABCDEF, 0x000000000000CDEF|
            \item \ \verb|0x0123456789ABCDEF, 0xFFFFFFFFFFFFCDEF|
            \item \ \verb|0x00000000FFFFCDEF, 0xFFFFFFFFFFFFCDEF|
            \item \ \verb|0x0123456789ABCDEF, 0xFFFFFFFFFFFFCDEF|
            \item \ \verb|0xFFFFFFFF89ABCDEF, 0xFFFFFFFFFFFFCDEF|
        \end{compactenum}
        \pro 下列操作不等价的是：
        \begin{choices}
            \item \verb|movzbq| 和 \verb|movzbl|
            \item \verb|movzwq| 和 \verb|movzwl|
            \item \verb|movl| 和 \verb|movslq|
            \item \verb|movslq %eax, %rax| 和 \verb|cltq|
        \end{choices}
        \sol C。A、B 中，\verb|movzbl/movzwl| 都生成了四字节，把高位设为0。D 中 \verb|cltq| 是对 \verb|%eax|的符号拓展。而 C 中 \verb|movl| 和 \verb|movzlq| 等价。
        \pro 判断下列 x86-64 ATT 数据传送指令是否合法。
            \qn \verb|movl $0x400010, $0x800010|
            \qn \verb|movl $0x400010, 0x800010|
            \qn \verb|movl 0x400010, 0x800010|
            \qn \verb|movq $-4, (%rsp)|
            \qn \verb|movq $0x123456789AB, %rax|
            \qn \verb|movabsq $0x123456789AB, %rdi|
            \qn \verb|movabsq $0x123456789AB, 16(%rcx)|
            \qn \verb|movq 8(%rsp), %rip|
        \sol 1 不正确，目的不能是立即数；2 正确；3 不正确，两个操作数不能同时是内存地址；4 正确；5 不正确，要用 \verb|movabsq|；6 正确；7 不正确，\verb|movabsq| 的目标地址必须是整数寄存器；8 不正确，不能用 \verb|mov| 向 \verb|%rip| 中传入数据，否则读程序计数器产生破坏。
        \pro 在 32 位机器中有如下定义 \texttt{int array[10] = \{0, 1, 2, 3, 4, 5, 6, 7, 8, 9\};}。
        
        某一时刻，\verb|%ecx| 存着第一个元素的地址，\verb|%ebx| 值为 3，那么下列操作中，哪一个将 \verb|array[3]| 移入了 \verb|%eax|？
        \begin{choices}
            \item \verb|leal 12(%ecx), %eax|
            \item \verb|leal (%ecx, %ebx, 4), %eax|
            \item \verb|movl (%ecx, %ebx, 4), %eax|
            \item \verb|movl 8(%ecx, %ebx, 2), %eax|
        \end{choices}
        \pro 下面的汇编代码对应一个 C 函数，其原型为 \verb|long func(long a, long b);|。将其翻译为 C 代码。
        \begin{minted}[frame=single, fontsize=\small]{gas}
    // a in %rdi, b in %rsi
    func:
        movq   %rdi, %rax
        salq   $4, %rax
        subq   %rdi, %rax
        movq   %rax, %rdi
        leaq   0(, %rsi, 8), %rax
        subq   %rsi, %rax
        addq   %rdi, %rax
        ret
        \end{minted}
        \sol 如下所示:
        \begin{minted}[frame=single, fontsize=\small]{c}
    long func(long a, long b) {
        return a * 15 + b * 7;
    }
        \end{minted}
        \pro 指令 \verb|setg %al| 会让寄存器 \verb|%al| 得到：
        \begin{choices}
            \item \verb+~(SF ^ OF) & ~ZF+
            \item \verb+~(SF | OF) & ~ZF+
            \item \verb+~(SF | OF)+
            \item \verb+~(SF ^ OF)+
        \end{choices}
        \sol A。如果 \verb|a > b|，那么 \verb|a – b| 或者为正，或者发生上溢变负；同时 \verb|a| 不等于 \verb|b|。
        \pro 下面的汇编代码对应一个 C 函数，其原型为 \verb|long func(long a, long b);|。将其翻译为 C 代码。
        \begin{minted}[frame=single, fontsize=\small]{gas}
    // a in %rdi, b in %rsi
    func: 
        movl    $1, %eax
        jmp     .L2
    .L4: 
        testb   $1, %sil
        je      .L3
        imulq   %rdi, %rax
    .L3: 
        sarq    %rsi
        imulq   %rdi, %rdi
    .L2: 
        testq   %rsi, %rsi
        jg      .L4
        rep ret 
        \end{minted}
        \sol 如下所示（快速幂）：
        \begin{minted}[frame=single, fontsize=\small]{c}
    long func(long a, long b) {
        long ans = 1;
        while (b > 0) {
            if (b & 1)
                ans = ans * a;
            b = b >> 1;
            a = a * a;
        }
        return ans;
    }
        \end{minted}
        \pro 对于下列四个函数，假设 gcc 开了编译优化，判断 gcc 是否会将其编译为条件传送。
        \qn
        \begin{minted}[frame=single, fontsize=\small]{c}
    long f1(long a, long b) { return (++a > --b) ? a : b; }
        \end{minted}
        \qn
        \begin{minted}[frame=single, fontsize=\small]{c}
    long f2(long a, long b) { return (*a > *b) ? --(*a) : (*b)--; }
        \end{minted}
        \qn
        \begin{minted}[frame=single, fontsize=\small]{c}
    long f3(long a, long b) { return a ? *a : (b ? *b : 0); }
        \end{minted}
        \qn
        \begin{minted}[frame=single, fontsize=\small]{c}
    long f4(long a, long b) { return (a > b) ? a++ : ++b; }
        \end{minted}
        \sol 1、4 会被编译成条件传送。1 由于比较前计算出的 \verb|a| 与 \verb|b| 就是传送的目标，因此会被编译成条件传送；2 由于比较结果会导致 \verb|a| 与 \verb|b| 指向的元素发生不同的改变，因此会被编译成条件跳转；3 由于指针 \verb|a| 可能无效，因此会被编译为条件跳转；4 会被编译成条件传送，因为 \verb|a| 和 \verb|b| 都是局部变量，返回的时候对 \verb|a| 和 \verb|b| 的操作都是无影响的。
        \pro 根据下面的汇编指令补充机器码中缺失的字节。
        \begin{table}[H]
            \centering
            \begin{tabular}{|c|c|}
                \hline
                机器码 & 汇编指令 \\ \hline
                \tt
                \begin{tabular}{l}
                    loop: \\
                    4004d0: 48 89 f8 \\
                    4004d3: eb \verb|_____| \\
                    4004d5: 48 d1 f8 \\
                    4004d8: 48 85 c0 \\
                    4004db: 7f \verb|_____| \\
                    4004dd: f3 c3 
                \end{tabular} &
                \begin{tabular}{l}
                    \\
                    \verb|mov   %rdi, %rax| \\
                    \verb|jmp   4004d8 <loop+0x8>| \\
                    \verb|sar   %rax| \\
                    \verb|test  %rax, %rax| \\
                    \verb|jg    4004d5 <loop+0x5>| \\
                    \verb|repz retq|
                \end{tabular}
                \\ \hline
            \end{tabular}
        \end{table}
        \sol \verb|03; f8|。第五行跳转位置为 \verb|0xf8(-8) + 0x4004dd = 0x4004d5|，注意进行跳转之前，PC 指向该指令的下一条指令。
        \pro 使用 GDB 查看某个可执行文件，发现其一段内存为：
        \begin{minted}[frame=single, fontsize=\small]{text}
    0x400598: 0x0000000000400488 0x0000000000400488
    0x4005a8: 0x000000000040048b 0x0000000000400493
    0x4005b8: 0x000000000040049a 0x0000000000400482
    0x4005c8: 0x000000000040049a 0x0000000000400498
        \end{minted}
        根据以下汇编代码，
        \begin{minted}[frame=single, fontsize=\small]{gas}
    0x400474: cmp   $0x7, %edi
    0x400477: ja    0x40049a
    0x400479: mov   %edi, %edi
    0x40047b: jmpq  *0x400598(, %rdi, 8)
    0x400482: mov   $0x15213, %eax
    0x400487: retq
    0x400488: sub   $0x5, %edx
    0x40048b: lea   0x0(, %rdx, 4), %eax
    0x400492: retq
    0x400493: mov   $0x2, %edx
    0x400498: and   %edx, %esi
    0x40049a: lea   0x4(%rsi), %eax
    0x40049d: retq             
        \end{minted}
        补全主函数的 C 代码。
        \begin{minted}[frame=single, fontsize=\small]{c}
    // a in %rdi, b in %rsi, c in %rdx
    int main(int a, int b, int c) {
        int res = 4;
        switch (a) {
            case 0:
            case 1:
                __________;
            case ___:
                res = __________;
                break;
            case ___:
                res = __________;
                break;
            case 3:
                __________;
            case 7:
                __________;
            default:
                __________;
        }
        return res;
    }
        \end{minted}
        \sol 做法是首先根据跳转表, 确定各标号的起始地址, 然后再作汇编代码的对应. 答案如下所示：
        \begin{minted}[frame=single, fontsize=\small]{c}
    case 0:
    case 1:
        c = c - 5;
    case 2:
        res = 4 * c; // or res *= c
        break;
    case 5:
        res = 86547; // or 0x15213
        break;
    case 3:
        c = 2;
    case 7:
        b = b & c;
    default:
        res += b; // or res = b + 4
        \end{minted}
        \pro 将下列汇编代码翻译成 C 代码。
        \begin{minted}[frame=single, fontsize=\small]{gas}
    func:
        movq   %rsi, %rax
        testq  %rdi, %rdi
        jne    .L7
        rep ret
    .L7:
        subq   $8, %rsp
        imulq  %rdi, %rax
        movq   %rax, %rsi
        subq   $1, %rdi
        call   func
        addq   $8, %rsp
        ret
        \end{minted}
        \begin{minted}[frame=single, fontsize=\small]{c}
    long func(long n, long m) {
        if (__________)
            return __________;
        return func(__________, __________);
    }
        \end{minted}
        \pro 将下列 C 代码翻译为汇编代码。
        \begin{minted}[frame=single, fontsize=\small]{c}
    void callee(long *a, long *b) {
        if (a == b) return;
        *a ^= *b;
        *b ^= *a;
        *a ^= *b;
    }
    void caller(long n, long arr[]) {
        for (long i = 0; i < n/2; i++)
            callee(&arr[i], &arr[n-i]);
    }
        \end{minted}
        汇编代码为:
        \begin{minted}[frame=single, fontsize=\small]{gas}
    callee:
        cmpq   %rsi, %rdi
        je     .L1
        movq   (%rsi), %rax
        xorq   (%rdi), %rax
        movq   ____________
        xorq   (%rsi), %rax
        movq   %rax, (%rsi)
        xorq   %rax, (%rdi)
    .L1:
        rep ret

    caller:
        pushq  %r12
        pushq  %rbp
        pushq  %rbx
        movq   %rdi, %rbp
        movq   %rsi, %r12
        movl   $0, %ebx
        jmp    .L4
    .L5:
        movq   %rbp, %rax
        subq   %rbx, %rax
        ____   (%r12, %rax, ____), %rsi
        ____   (%r12, %rbx, ____), %rdi
        call   callee
        addq   $1, %rbx
    .L4:
        movq   %rbp, %rax
        shrq   $63, %rax
        addq   %rbp, %rax
        sarq   %rax
        cmpq   %rbx, %rax
        jg     .L5
        popq   ____
        popq   ____
        popq   ____
        ret
        \end{minted}
        
        对于上面代码，在 x86-64、操作系统为 Linux 的情况下，假设 main 在 \verb|0x4000ac| 处调用 caller， caller 在 \verb|0x400088| 处调用 callee；调用函数（\verb|call xx|）的代码长度为 5。在 main 即将调用 caller 时，部分寄存器的情况见下表左侧。请在下图右侧画出控制流第一次走 到 \verb|.L1| 时，堆栈的结构。
        \begin{table}[H]
            \centering
            \begin{tabular}{|c|c|c|c|}
                \hline
                寄存器 & 调用前的值 & 地址 & 内容（不确定的空格不用填） \\ \hline
                \verb|%rsp| & \verb|0xffffffffff80| & \verb|0xf...f88~8f| &  \\ \hline
                \verb|%rax| & \verb|0x0| & \verb|0xf...f80~87| &  \\ \hline
                \verb|%rbx| & \verb|0x15| & \verb|0xf...f78~7f| &  \\ \hline
                \verb|%rbp| & \verb|0x18| & \verb|0xf...f70~77| &  \\ \hline
                \verb|%r12| & \verb|0x213| & \verb|0xf...f68~6f| &  \\ \hline
                \verb|%rsi| & \verb|0x0| & \verb|0xf...f60~67| &  \\ \hline
                \verb|%rdi| & \verb|0x0| & \verb|0xf...f58~5f| &  \\ \hline
                &  & \verb|0xf...f50~57| &  \\ \hline
            \end{tabular}
        \end{table}
        \sol 阅读 \verb|callee| 汇编代码可知，首先 \verb|*b| 进入 \verb|%rax|，随后 \verb|%rax| 变为 \verb|*a ^ *b|。根据 C 代码确定下一步要将结果转入 \verb|*a| 中，所以汇编语言代码的空应填写 \verb|%rax, %(rdi)|。\verb|caller| 汇编代码的阅读不需要太仔细(就做题而言)，极易填写。注意这里体现的典型的 \verb|for| 循环结构，快速识别 \verb|%ebx| 为循环变量，\verb|.L4| 后为循环条件检查。后者需要计算 \verb|n / 2|，要留意编译器选择用移位作除法，但是对于负数需要作修正(舍入方式不同)。\verb|sarq| 等移位操作若只有一个操作数，则移位量默认为 1。

        此处栈空间的填写难度不大。首先，\verb|main| 即将发生调用时，\verb|%rsp| 值为 \verb|0xf...f80|，接下来 \verb|call| 后需要在栈中设置返回地址，这个返回地址为 \verb|0x4000ac + 0x5 = 0x4000b1|。所以栈的第三个空填此。注意栈是向下增长的，所以前两个空不确定。接下来 \verb|caller| 要保存 \verb|%r12, %rbp, %rbx| 三个寄存器，因此先后填入 \verb|0x213, 0x18, 0x15|。最后，\verb|callee| 被调用，需要再次填入返回地址 \verb|0x400088 + 0x5 = 0x40008d|。最后一个空显然不确定。
        \pro 在 x86-64、Linux 操作系统下有如下 C 定义：
        \begin{minted}[frame=single, fontsize=\small]{c}
    struct A {
        char CC1[6];
        int II1;
        long LL1;
        char CC2[10];
        long LL2;
        int II2;
    };
        \end{minted}
            \qn \verb|sizeof(A) = |\rule{2.5cm}{0.25mm}。
            \qn 将 \texttt{A} 重排后，令结构体尽可能小，那么得到的新的结构体大小为 \rule{2.5cm}{0.25mm} 字节。
        \pro 在 x86-64、LINUX 操作系统下，考虑如下的 C 定义：
        \begin{minted}[frame=single, fontsize=\small]{c}
    typedef union {
        char c[7];
        short h;
    } union_e;

    typedef struct {
        char d[3];
        union_e u;
        int i;
    } struct_e;

    struct_e s;
        \end{minted}
            \qn \verb|s.u.c| 的首地址相对于 \verb|s| 的首地址的偏移量是 \rule{2.5cm}{0.25mm} 字节。
            \qn \verb|sizeof(union_e) = |\rule{2.5cm}{0.25mm} 字节。
            \qn \verb|s.i| 的首地址相对于 \verb|s| 的首地址的偏移量是 \rule{2.5cm}{0.25mm} 字节。
            \qn \verb|sizeof(struct_e) = |\rule{2.5cm}{0.25mm} 字节。
            \qn 若只将 \verb|i| 的类型改成 \verb|short|，那么 \verb|sizeof(struct_e) = |\rule{2.5cm}{0.25mm} 字节。
            \qn 若只将 \verb|h| 的类型改成 \verb|int|，那么 \verb|sizeof(union_e) = |\rule{2.5cm}{0.25mm} 字节。
            \qn 若将 \verb|i| 的类型改成 \verb|short|、\verb|h| 的类型改成 \verb|int|，则 \verb|sizeof(union_e) = |\rule{2.5cm}{0.25mm} 字节，\verb|sizeof(struct_e) = |\rule{2.5cm}{0.25mm} 字节。
            \qn 若只将 \verb|short h| 的定义删除，那么 (1)\textasciitilde(4) 问的答案分别是 \rule{1cm}{0.25mm}、\rule{1cm}{0.25mm}、\rule{1cm}{0.25mm}、\rule{1cm}{0.25mm} 字节。
        \pro 以下提供了一段代码的 C 语言、汇编语言以及运行到某一时刻栈的情况。
        \begin{minted}[frame=single, fontsize=\small]{gas}
0000000000400596 <func>: 
    400596: sub    $0x28, %rsp
    40059a: mov    %fs:0x28, %rax
    4005a3: mov    %rax, 0x18(%rsp)
    4005a8: xor    %eax, %eax
    4005aa: mov    (%rdi), %rax
    4005ad: mov    0x8(%rdi), %rdx
    4005b1: cmp    %rdx, %rax
    4005b4: jge    _____(1)_____
    4005b6: mov    %rdx, (%rdi)
    4005b9: mov    %rax, 0x8(%rdi)
    4005bd: mov    0x8(%rdi), %rax
    4005c1: test   %rax, %rax
    4005c4: jne    4005cb <func+0x35>
    4005c6: mov    (%rdi), %rax
    4005c9: jmp    _____(2)_____
    4005cb: mov    (%rdi), %rdx
    4005ce: sub    %rax, %rdx
    4005d1: mov    %rdx, (%rsp)
    4005d5: mov    %rax, 0x8(%rsp)
    4005da: mov    _____(3)_____, %rdi
    4005dd: callq  400596 <func>
    4005e2: mov    0x18(%rsp), %rcx
    4005e7: xor    _____(4)_____, %rcx
    4005f0: _____(5)_____    4005f7 <func+0x61>
    4005f2: callq  400460 <__stack_chk_fail@plt>
    4005f7: add    _____(6)_____, %rsp
    4005fb: retq
 
00000000004005fc <main>: 
    4005fc: sub    $0x28, %rsp 
    400600: mov    %fs:0x28, %rax 
    400609: mov    %rax, 0x18(%rsp) 
    40060e: xor    %eax, %eax 
    400610: movq   0x69, (%rsp) 
    400618: movq   0xfc, 0x8(%rsp) 
    400621: mov    %rsp, %rdi 
    400624: callq  400596 <func>
    400629: mov    %rax, %rsi
    40062c: mov    $0x4006e4, %edi
    400631: mov    $0x0, %eax
    400636: callq  400470 <printf@plt>
    40063b: mov    0x18(%rsp), %rdx
    400640: xor    _____(4)_____, %rdx
    400649: _____(5)_____    400650 <main+0x54>
    40064b: callq  400460 <__stack_chk_fail@plt>
    400650: mov    $0x0, %eax
    400655: add    _____(6)_____, %rsp
    400659: retq
        \end{minted}
        C 语言代码：
        \begin{minted}[frame=single, fontsize=\small]{c}
    typedef struct{
        long a;
        long b;
    } pair_type;

    long func(pair_type *p) {
        if (p -> a < p -> b) {
            long temp = p -> a;
            p -> a = p -> b;
            p->b = temp;
        }
        if (_____(7)_____) {
            return p->a;
        }
        pair_type np;
        np.a = _____(8)_____;
        np.b = _____(9)_____;
        return func(&np);
    } 

    int main(int argc, char* argv[]) {
        pair_type np;
        np.a = _____(10)_____;
        np.b = _____(11)_____;
        printf("%ld", func(&np));
    }
        \end{minted}
        堆栈情况如下（从高地址向低地址列举）：
        {\tt \begin{compactenum}
            \item\ 0x0000000000000000
            \item\ 0xc76d5add7bbeaa00 
            \item\ 0x00007fffffffdf60 
            \item\ \verb|__________________|
            \item\ \verb|__________________|
            \item\ 0x0000000000400629 
            \item\ \verb|__________________|
            \item\ \verb|__________________|
            \item\ 0x0000000000000001 
            \item\ 0x0000000000000069 
            \item\ 0x0000000000000093 
            \item\ \verb|__________________|
            \item\ 0x00000000ff000000 
            \item\ \verb|__________________|
            \item\ 0x0000000000000000 
            \item\ \verb|__________________|
            \item\ \verb|__________________|
            \item\ \verb|__________________|
            \item\ 0x0000000000000000 
            \item\ \verb|__________________|
            \item\ \verb|__________________|
            \item\ 0x000000000000002a 
            \item\ 0x000000000000003f 
            \item\ 0x00000000004005e2 
        \end{compactenum}}
        一些可能用到的字符的 ASCII 码如下：
        \begin{table}[H]
            \centering
            \begin{tabular}{|c|c|c|c|c|c|c|c|c|c|}
                \hline
                换行 & 空格 & \verb|"| & \verb|%| & \verb|(| & \verb|)| & \verb|,| & \verb|0| & \verb|A| & \verb|a| \\ \hline
                \verb|0x0a| & \verb|0x20| & \verb|0x22| & \verb|0x25| & \verb|0x28| & \verb|0x29| & \verb|0x2c| & \verb|0x30| & \verb|0x41| & \verb|0x61| \\ \hline
            \end{tabular}
        \end{table}

        回答下列问题。
            \qn gdb 下使用命令 \verb|x/4b 0x4006e4| 后（即查看 \verb|0x4006e4| 开始的 4 个字节，用 16 进制表示）得到的输出结果是 \rule{2.5cm}{0.25mm}。
            \qn 互相翻译 C 语言代码和汇编代码，补充缺失的空格（标号相同的为同一格）。
            \qn 补充栈的内容。使用16进制，可以不写前导多余的 0；对于给定已知条件后仍无法确定的值，填写“不确定”；已知程序运行过程中寄存器 \verb|%fs| 的值没有改变。
            \qn 程序运行结果是 \rule{2.5cm}{0.25mm}。
        \sol 这是典型的汇编综合题，需要同时结合 C 代码和汇编代码来观察。
        
        第一步先通读 C 代码，确定代码大意，便于在汇编中找一些对应的片段。然后来填写汇编代码。主要难点在厘清 C 语言变量，内存和寄存器的映射关系，确定代码的用途。先看主函数，前三行为金丝雀。接下来，\verb|0x69, 0xfc| 分别进入 \verb|%rsp, %rsp + 8| 对应的内存位置中。根据主函数唯一的 \verb|&np| 取局部变量地址，以及紧接着 \verb|%rsp| 被作为参数传入 \verb|func|，知道这就是 \verb|np| 结构体的初始化工作。所以 C 代码的第 10、11 空分别填写 105、252。

        接着，\verb|func| 的返回值被存入 \verb|%rsi|, 一个立即数（看上去像地址）被存入 \verb|%rdi|，然后调用 \verb|printf| 函数。结合 C 代码知道 \verb|%rdi| 是字符串 \verb|"%ld"|（属于 \verb|char*| 类型）的地址。所以由给出的 ASCII 码表计算得到问答题(1)应填写 \verb|0x25, 0x6c, 0x64, 0x00|。特别留意字符串的最后有一个 \verb|0x00| 表示结尾！主函数的最后几个空是检查金丝雀是否被破坏，易知金丝雀的原值在 \verb|%fs:0x28|，而在栈上对应于地址 \verb|%rsp + 0x18|（这个之后还要用到）。检查金丝雀是否被破坏是容易的，\verb|40063b| 先把后者放到 \verb|%rcx| 中，所以要比较 \verb|%fs:0x28, %rdx|，这就是空(4)的答案。接下来要条件跳转，因为如果条件成立会正常返回，所以空(5)写 \verb|je|。空(6)是函数返回释放栈空间，因为一开始分配了 \verb|0x28| 个字节，所以这个空也要填 \verb|0x28|。

        现在分析 \verb|func| 函数。前三行是金丝雀以及寄存器清零。从地址 \verb|4005aa| 开始分析，首先 \verb|p->a, p->b| 分别进入 \verb|%rax, %rdx| 中。接下来比较 \verb|p->a >= p->b|。结合 C 代码，如果这成立，就不需要经过交换 \verb|p->a, p->b| 的过程。立刻看出，\verb|4005b6, 4005b9| 两行完成了交换工作（没有出现 \verb|temp| 的对应寄存器映射）。由于不确定交换之后 \verb|%rax| 到底对应 \verb|p->a, p->b| 中的哪个，所以其要重新移入 \verb|p->b| 到 \verb|p->a| 中。所以 \verb|4005bd| 就是空(1)的跳转目标。接下来用 \verb|test| 指令判断 \verb|%rax| 是否为 0。注意 C 代码，如果某个条件成立，就要 \verb|return p->a|。因为 \verb|4005cb| 把 \verb|p->a| 放到 \verb|%rax| 中，所以空(2)会跳到函数结尾返回，也就是说跳转目标是 \verb|4005e2|（返回之前要检查金丝雀是否被破坏）。由此也可知只有 \verb|p->b| 为零时才返回，空(7)填写 \verb|p->b == 0|。

        紧接着处理 \verb|func| 的递归调用。结合 C 代码知道要生成新的 \verb|np|，而且 \verb|np.a, np.b| 分别在栈地址的 \verb|%rsp, %rsp + 8| 处。因为这时 \verb|p->a| 在 \verb|%rdx| 中，\verb|p->b| 在 \verb|%rax| 中，所以由 \verb|4005ce, 4005d1| 两行立刻知道(8)、(9)两空分别写 \verb|p->a - p->b, p->b|。

        补全 C 代码后，可以看出这是辗转相除法，输入是 252、105，它们的最大公约数是 21，故问答题的(4)回答 21。

        最后填写栈。对于递归调用，我们应该先正确划分各次调用的栈的边界。第一步，由于主函数和 \verb|func| 一开始都 \verb|subq 0x28, %rsp|，所以栈空间都有 40 字节或 5 个 64 位地址。注意 6 号位置是 \verb|0x400629|，恰好是主函数调用 \verb|func| 的返回地址，所以我们知道 1--5 号位置是主函数的栈帧，6 是从 \verb|func| 返回主函数的返回地址；7--11 是第一次调用 \verb|func| 的栈帧，12 是从第二次递归调用返回第一次调用 \verb|func| 的返回地址，13--17 是第二次调用 \verb|func| 的栈帧，以此类推。所以 12、18 都填 \verb|0x4005e2|（\verb|func| 中递归调用的下一条指令的位置）。
        
        第二步，根据前面的分析，每个栈的栈帧的底部向上数第一个和第二个数都是存储局部变量 \verb|np|（无论是主函数还是 \verb|func|），所以位置 4、5 似乎分别应该写上 \verb|0xfc, 0x69|。但是，注意后面的递归调用中，这两个位置的数会被交换以保证 \verb|p->a >= p->b|！所以这两个空要倒过来，4、5 分别写上 \verb|0x69, 0xfc|。第三步，由前面的分析，金丝雀值一直存在 \verb|%rsp + 0x18|，也就是栈帧自底向上第四个位置，又题目条件说 \verb|%fs| 寄存器没有改变，所以所有的金丝雀值都是 \verb|0xc76d5add7bbeaa00|，这样 8、14、20 都是这个值。

        对于栈的 10、11 号位置，前面的分析知道它们应该是 \verb|252 - 105, 105 = 0x93, 0x69|（顺序未定），而交换后保证上一个要小于下一个，所以 10、11 给出的情况印证我们的分析。同理，16、17 空是 \verb|0x93 - 0x69 = 0x2a, 0x69|，上一个要小于下一个。虽然 22、23 号位置不用填写，但是我们务必注意，栈在这个状态时，刚刚进行下一次递归调用，这两个数还没有被下一次递归调用修改顺序（尽管实际上不需要修改顺序）。

        栈的剩下位置，7、21 都是不确定的，因为只是分配了没有修改过。（注意栈帧开头应该不是旧的帧指针，因为如果要这样需要显式 \verb|push|。）
    \end{problems}

\chapter{汇编语言{---}往年考题}
    \begin{problems}
        \proy{2018} 在 x86-64 下，以下哪个选项的说法是错误的？
        \begin{choices}
            \item \verb|movl| 指令以寄存器作为目的时，会将该寄存器的高位 4 字节设置为 0
            \item \verb|cltq| 指令的作用是将 \verb|%eax| 符号扩展到 \verb|%rax|
            \item \verb|movabsq| 指令只能以寄存器作为目的
            \item \verb|movswq| 指令的作用是将零扩展的字传送到四字节目的
        \end{choices}
        \proy{2018} 下列关于程序控制结构的机器代码实现的说法中，正确的是：
        \begin{choices}
            \item 使用条件跳转语句实现的程序片段比使用条件赋值语句实现的同一程序片段的运行效率高
            \item 使用条件跳转语句实现的程序片段与使用条件赋值语句实现的同一程序片段虽然效率可能不同，但在 C 语言的层面上看总是有着相同的行为
            \item 一些 \verb|switch| 语句不会被 gcc 用跳转表的方式实现
            \item 以上说法都不正确
        \end{choices}
        \proy{2018} 下列关于条件码的叙述中，不正确的是：
        \begin{choices}
            \item 所有算术指令都会改变条件码
            \item 所有比较指令都会改变条件码
            \item 所有与数据传送有关的指令都会改变条件码
            \item 条件码一般不会直接读取，但可以直接修改
        \end{choices}
        \sol 此题有错，A、B 都不准确，有一些 SIMD 指令并不会改变条件码；C 明显是错的；D 基本正确，确实有指令可以直接写条件码。
        \proy{2017} 在下列的 x86-64 汇编代码中，错误的是：
        \begin{choices}
            \item \verb|movq %rax, (%rsp)|
            \item \verb|movl $0xFF, (%ebx)|
            \item \verb|movsbl (%rdi), %eax|
            \item \verb|leaq (%rdx, 1), %rdx|
        \end{choices}
        \sol 此题有错，B、D 都不对，其中 B 不能用 32 位寄存器寻址，D 中少了一个逗号。
        \proy{2017} 在下列关于条件传送的说法中，正确的是：
        \begin{choices}
            \item 条件传送可以用来传送字节、字、双字和四字的数据
            \item C语言中的“\verb|?:|”条件表达式都可以编译成条件传送
            \item 使用条件传送总可以提高代码的执行效率
            \item 条件传送指令不需要用后缀(例如 \verb|b, w, l, q|)来表明操作数的长度
        \end{choices}
        \proy{2017} 在下列指令中，其执行会影响条件码中的 CF 位的是：
        \begin{choices}
            \item \verb|jmp NEXT|
            \item \verb|jc NEXT|
            \item \verb|inc %bx|
            \item \verb|shl $1, %ax|
        \end{choices}
        \proy{2016} 下列关于比较指令 \verb|cmp| 说法中，正确的是：
        \begin{choices}
            \item 专用于有符号数比较
            \item 专用于无符号数比较
            \item 专用于串比较
            \item 不区分比较的对象是有符号数还是无符号数
        \end{choices}
        \proy{2016} 在如下代码段的跳转指令中，目的地址是：
        \begin{minted}[frame=single, fontsize=\small]{gas}
    400020: 74 F0 je   ________
    400022: 5d    pop  %rbp
        \end{minted}
        \begin{choices}
            \item \texttt{400010}
            \item \texttt{400012}
            \item \texttt{400110}
            \item \texttt{400112}
        \end{choices}
        \proy{2016} 对于如下的 C 语言中的条件转移指令，它所对应的汇编代码中至少包含几条条件转移指令：\verb+if (a > 0 && a != 1 || a < 0 && a != -1) b = a;+？
        \begin{choices}
            \item 2 条
            \item 3 条
            \item 4 条
            \item 5 条
        \end{choices}
        \sol B，这个条件相当于 \verb+(a == 0 || a == 1 || a == -1)+。
        \proy{2016} 将 \verb|%ax| 清零，下列指令不能实现该效果的是：
        \begin{choices}
            \item \verb|sub %ax, %ax|
            \item \verb|xor %ax, %ax|
            \item \verb|test %ax, %ax|
            \item \verb|and $0, %ax|
        \end{choices}
        \proy{2016} 在如下 \verb|switch| 语句翻译得到的跳转表中，哪些标号没有出现在分支中？
        \begin{minted}[frame=single, fontsize=\small]{gas}
    addq    $1, %rdi
    cmpq    $8, %rdi
    ja      .L2
    jmp     *.L4(, %rdi, 8)
    .L4:
        .quad .L9
        .quad .L5
        .quad .L6
        .quad. L7
        .quad .L2
        .quad .L7
        .quad .L8
        .quad .L2
        .quad .L5
        \end{minted}
        \begin{choices}
            \item 3, 6
            \item $-1$, 4
            \item 0, 7
            \item 2, 4
        \end{choices}
        \proy{2016} 已知短整型数组 \verb|S| 的起始地址和下标 \verb|i| 分别存放在寄存器 \verb|%rdx| 和 \verb|%rcx|，将 \verb|&S[i]| 存放在寄存器 \verb|%rax| 中所对应的汇编代码是：
        \begin{choices}
            \item \verb|leaq (%rdx, %rcx, 1), %rax|
            \item \verb|movw (%rdx, %rcx, 2), %rax|
            \item \verb|leaq (%rdx, %rcx, 2), %rax|
            \item \verb|movw (%rdx, %rcx, 1), %rax|
        \end{choices}
        \proy{2015} 下列寻址模式中，正确的是：
        \begin{choices}
            \item \verb|(%eax, , 4)|
            \item \verb|(%eax, %esp, 3)|
            \item \verb|123|
            \item \verb|$1(%ebx, %ebp, 1)|
        \end{choices}
        \proy{2015} 假设某条 C 语言 \verb|switch| 语句编译后产生了如下的汇编代码及跳转表：
        \begin{minted}[frame=single, fontsize=\small]{gas}
    movl   8(%ebp), %eax
    subl   $48, %eax
    cmpl   $8, %eax
    ja     .L2
    jmp    *.L7(, %eax, 4)
    .L7:
        .long .L3
        .long .L2
        .long .L2
        .long .L5
        .long .L4
        .long .L5
        .long .L6
        .long .L2
        .long .L3
        \end{minted}
        在源程序中，下面的哪些标号出现过？
        \begin{choices}
            \item \verb|'2', '7'|
            \item \verb|1|
            \item \verb|'3'|
            \item \verb|5|
        \end{choices}
        \sol C，从 \verb|subl| 可看出待比较的数是字符型。
        \proy{2015, 2014} 下列的指令组中，哪一组指令只改变条件码，而不改变寄存器的值？
        \begin{choices}
            \item \verb|CMP, SUB|
            \item \verb|TEST, AND|
            \item \verb|CMP, TEST|
            \item \verb|LEAL, CMP|
        \end{choices}
        \proy{2014} 简单的 \verb|switch| 语句常采用跳转表的方式实现，在 x86-64 系统中，下述最有可能是正确的 \verb|switch| 分支跳转汇编指令的是哪个？
        \begin{choices}
            \item \verb|jmp .L3(, %eax, 4)|
            \item \verb|jmp .L3(, %eax, 8)|
            \item \verb|jmp *.L3(, %eax, 4)|
            \item \verb|jmp *.L3(, %eax, 8)|
        \end{choices}
        \proy{2018} 以下代码的输出结果不可能是
        \begin{minted}[frame=single, fontsize=\small]{c}
    union {
        double d;
        struct {
            int i;
            char c[4];
        } s;
    } u;
    u.d = 1;
    printf("%d\n", u.s.c[2]);
        \end{minted}
        \begin{choices}
            \item 0
            \item $-16$
            \item 240
            \item 191
        \end{choices}
        \sol D。题目没有说明大小端，如果是大端法则为 A。若为小端法，注意 \verb|char| 在不同的编译器上可能是有符号型或者是无符号型，所以计算知 $-16$ 和 240 都有可能。
        \proy{2018} 下列关于 C 语言中的结构体（struct）以及联合（union）的说法中，正确的是：
        \begin{choices}
            \item 对于任意 struct，将其成员按照其实际占用内存大小从小到大的顺序进行排列不一定会使之内存占用最小
            \item 对于任意 struct，将其成员按照其实际占用内存大小从小到大的顺序进行排列一定不会使之内存占用最大
            \item 对于任意 union，将其成员按照其实际占用内存大小从小到大的顺序进行排列不一定会使之内存占用最小
            \item 对于任意 union，将其成员按照其实际占用内存大小从小到大的顺序进行排列一定不会使之内存占用最大
        \end{choices}
        \proy{2017} 有如下代码段：
        \begin{minted}[frame=single, fontsize=\small]{c}
    int func(int x, int y);
    int (*p) (int a, int b);
    p = func;
    p(0, 0);
        \end{minted}
        对应的下列 x86-64 过程调用正确的是：
        \begin{choices}
            \item \verb|call *%rax|
            \item \verb|call *(%rax)|
            \item \verb|call (%rax)|
            \item \verb|call func|
        \end{choices}
        \sol 此题有错，A、D 都可能产生（和优化选项有关）。
        \proy{2017} 有定义：\verb|int A[3][2] = {{1,2}，{3,3}，{2,1}};|，则 \verb|A[2]| 是：
        \begin{choices}
            \item \verb|&A + 16|
            \item \verb|A + 16|
            \item \verb|*A + 4|
            \item \verb|*A + 2|
        \end{choices}
        \proy{2015} 已知下面的数据结构，假设在 Linux/IA32 下要求对齐，这个结构的总的大小是多少个字节？如果重新排列其中的字段，最少可以达到多少个字节？
        \begin{minted}[frame=single, fontsize=\small]{c}
    struct {
        char a;
        double *b;
        double c;
        short d;
        long long e;
        short f;
    };
        \end{minted}
        \begin{choices}
            \item 32, 28
            \item 36, 32
            \item 28, 26
            \item 26, 26
        \end{choices}
        \proy{2015} 在“大端法”下，已知如下的 C 语言数据结构：\verb|union { char c[2]; int i; };|。当 \verb|c| 的值为 \verb|0x01, 0x23| 时，\verb|i| 的值为：
        \begin{choices}
            \item \verb|0x0123|
            \item \verb|0x2301|
            \item \verb|0x01230000|
            \item 不确定
        \end{choices}
        \proy{2014} 有如下定义的结构，在 x86-64 下，下述结论中错误的是？
        \begin{minted}[frame=single, fontsize=\small]{c}
    struct {
        char c;
        union {
            char vc;
            double value;
            int vi;
        } u;
        int i;
    } sa;
        \end{minted}
        \begin{choices}
            \item \verb|sizeof(sa) == 24|
            \item \verb|(&sa.i - &sa.u.vi) == 8|
            \item \verb|(&sa.u.vc - &sa.c) == 8|
            \item 优化成员变量的顺序，可以做到 \verb|sizeof(sa) == 16|
        \end{choices}
        \proy{2013} 32 位 x86、Windows 操作系统下定义 struct \verb|S| 包含: \verb|double a, int b, char c|, 请问 \verb|S| 在内存空间中最多和最少分别能占据多少个字节（32 位 Windows 系统按 1、4、8 的原则对齐 \verb|char, int, double|）？
        \begin{choices}
            \item 16, 13
            \item 16, 16
            \item 24, 13
            \item 24, 16
        \end{choices}
    \end{problems}
	\chapter{Bomblab 回顾}
    \begin{summary}
        \begin{compactitem}
            \item 熟练掌握汇编语言语法，会阅读汇编语言中的各种程序结构，能和 C 语言进行快速转换。
            \item 知道如何使用 gdb 进行调试。
        \end{compactitem}
    \end{summary}

    \begin{problems}
        \pro 写出使用 gcc 编译源代码 \verb|lab.c|、生成可执行文件 \verb|lab|、采用二级编译优化的命令。
        \sol \verb|gcc lab.c -o lab -O2|
        \pro 写出使用 gcc 编译源代码 \verb|foo.c|、生成汇编语言文件 \verb|foo.s| 的命令。
        \sol \verb|gcc foo.c -S|
        \pro 将可执行文件逆向工程为汇编代码的工具是：
        \begin{choices}
            \item gcc
            \item gdb
            \item objdump
            \item hexedit
            \item gedit
        \end{choices}
        \sol C。其中 D 是编辑二进制文件的工具，E 是一种图形化的编辑器。
        \pro Gdb 中，单步指令执行的命令是：
        \begin{choices}
            \item \verb|r|
            \item \verb|b|
            \item \verb|p|
            \item \verb|finish|
            \item \verb|si|
            \item \verb|disas|
        \end{choices}
        \sol E。A 为运行，B 为设置断点，C 为打印某些信息，D 为停止调试，E 为显示反汇编代码。
        \proy{2013} 在完成本课程的 bomblab 的时候，通常先执行 \verb|gdb bomb| 启动调试，然后执行 \verb|___ explode_bomb| 命令以防引爆炸弹，之后在进行其他必要的设置后，最后执行 \verb|___| 命令以便开始执行程序。上述两个空格对应的命令是：
        \begin{choices}
            \item \verb|st, ru|
            \item \verb|br, go|
            \item \verb|br, ru|
            \item \verb|st, go|
        \end{choices}
        \sol C。
    \end{problems}
	\chapter{体系结构初步}
    \begin{summary}
        \begin{compactitem}
            \item 了解体系结构的简要发展史，熟悉 RISC、CISC 的区别和应用。了解基本的门电路、算术运算器、选择器、触发器和寄存器的原理和实现。
            \item 熟悉 Y86-64 “体系结构”中，各种指令的编码规则（包括操作数、指令子类型等），会将汇编代码和指令编码相互转换。简单了解 MIPS 体系结构。
            \item 在上一项的基础上，理解并熟练\uwave{记忆} Y86-64 处理器执行各种指令的几个阶段及其详细的功能（含重要的中间状态变化）。理解并熟练\uwave{记忆} Y86-64 顺序处理器的实现。
            \item 知道流水线设计的通用原理，会进行相关的计算和分析。
        \end{compactitem}
    \end{summary}

    \begin{problems}
        \pro 下列描述更符合（早期）RISC 还是 CISC？
            \qn 指令机器码长度固定。
            \qn 指令类型多、功能丰富。
            \qn 不采用条件码。
            \qn 实现同一功能，需要的汇编代码较多。
            \qn 译码电路复杂。
            \qn 访存模式多样。
            \qn 参数、返回地址都使用寄存器进行保存。
            \qn x86-64。
            \qn MIPS。
            \qn 广泛用于嵌入式系统。
            \qn 已知某个体系结构使用 \verb|add R1, R2, R3| 来完成加法运算。当要将数据从寄存器 \verb|S| 移动至寄存器 \verb|D| 时，需要使用 \verb|add S, #ZR, D| 进行操作（\verb|#ZR| 是一个恒为 0 的寄存器）。
            \qn 已知某个体系结构提供了 \verb|xlat| 指令，它以一个固定的寄存器 \verb|A| 为基地址，以另一个固定的寄存器 \verb|B| 为偏移量，在 \verb|A| 对应的数组中取出下标为 \verb|B| 的项的内容，放回寄存器 \verb|A| 中。
        \sol 依次为 RISC、CISC、RISC、RISC、CISC、CISC、RISC、CISC、RISC、RISC、RISC、CISC。这里对于 1，固定长的指令译码电路简单；通常 CISC 有变长指令，例如 x86 中清零用 \verb|xor|，因为这样可以节省指令长度。对于 10，嵌入式系统是指可穿戴设备、物联网之类的小设备，功能简单且需要省电，因此处理器电路不用太复杂。
        \pro 写出下列电路对应的逻辑表达式：
        \begin{figure}[H]
            \centering
            \begin{tikzpicture}[
                GateCfg/.style={
                    logic gate inputs={normal, normal, normal},
                    draw
                }
            ]
                \path (0,0) node[and gate US, GateCfg] (AND1) {} ++(2,-2) node[and gate US, GateCfg] (AND2) {} ++(5,1) node[or gate US, GateCfg] (OR1) {} (AND1.input 3) ++(-1,0) node[not gate US, draw] (N1) {} (AND2.input 3) ++(-1,0) node[not gate US, draw] (N2) {} (AND2.input 1 -| N1) node[not gate US, draw] (N3) {};
        
                \draw (OR1.input 1) -- ++(-1.5,0) |- (AND1.output) (OR1.input 3) -- ++(-1.5,0) |- (AND2.output) (N2.output) -- (AND2.input 3) (N1.output) -- (AND1.input 3) (N3.output) -- (AND2.input 1) (AND1.input 1) -- ++(-3,0) coordinate (init) node[anchor=east] {\texttt{A}} node[pos=0.6] (temp) {} (N1-| temp) ++(0,5pt) edge (temp.center) arc (90:-90:5pt) |- (N3.input) (init |- N1) node[anchor=east] {\texttt{B}} -- (N1.input) node[pos=0.4] (temp2) {} (temp2.center) |- (N2.input) (OR1.output) -- ++(2,0) node [midway,anchor=south]{\texttt{f(A, B)}};
            \end{tikzpicture}
        \end{figure}
        \sol \verb+f(A, B) = (!A && B) || (A && !B)+。
        \pro 根据 Y86-64 体系结构完成下表中的更新逻辑：
        \begin{table}[H]
            \tt
            \centering
            \begin{tabular}{|c|c|c|c|}
                \hline
                &  & call & jXX \\ \hline
                \multirow{4}{*}{Fetch (F)} & icode, ifun & {\qquad \qquad \qquad \qquad} & {\qquad \qquad \qquad \qquad} \\ \cline{2-4} 
                & rA, rB &  &  \\ \cline{2-4} 
                & valC &  &  \\ \cline{2-4} 
                & valP &  &  \\ \hline
                \multirow{2}{*}{Decode (D)} & valA, srcA &  &  \\ \cline{2-4} 
                & valB, srcB &  &  \\ \hline
                \multirow{2}{*}{Execute (E)} & valE &  &  \\ \cline{2-4} 
                & Cond Code &  &  \\ \hline
                Memory (M) & valM &  &  \\ \hline
                \multirow{2}{*}{Writeback (W)} & dstE &  &  \\ \cline{2-4} 
                & dstM &  &  \\ \hline
                PC & PC &  &  \\ \hline
            \end{tabular}
        \end{table}
        \sol 首先确定 \verb|call V| 和 \verb|jXX V| 两个指令都长 9 字节，然后按运算逻辑进行填写即可。没有填写的空格表示不需要填。
        \begin{table}[H]
            \tt
            \centering
            \begin{tabular}{|c|c|c|c|}
                \hline
                &  & call & jXX \\ \hline
                \multirow{4}{*}{Fetch} & icode:ifun & icode:ifun <- M1[PC] & icode:ifun <- M1[PC] \\ \cline{2-4} 
                & rA, rB &  &  \\ \cline{2-4} 
                & valC & valC <- M8[PC+1] & valC <- M8[PC+1] \\ \cline{2-4} 
                & valP & valP <- PC+9 & valP <- PC+9 \\ \hline
                \multirow{2}{*}{Decode} & valA &  &  \\ \cline{2-4} 
                & valB & valB <- R[\%rsp] &  \\ \hline
                \multirow{2}{*}{Execute} & valE & valE <- valB+(-8) &  \\ \cline{2-4} 
                & Cond Code &  & Cnd <- Cond(CC, ifun) \\ \hline
                Memory (M) & valM & M8[valE] <- valP &  \\ \hline
                \multirow{2}{*}{Writeback} & dstE & R[\%rsp] <- valE &  \\ \cline{2-4} 
                & dstM &  &  \\ \hline
                PC & PC & PC <- valC & PC <- Cnd?valC:valP \\ \hline
            \end{tabular}
        \end{table}
        \pro 已知 \verb|valC| 为指令中的常数值，\verb|valM| 为访存得到的数据，\verb|valP| 为 \verb|PC| 自增得到的值，完成以下的 PC 更新逻辑：
        \begin{minted}[frame=single, fontsize=\small]{text}
    int new_pc = [
        icode == ICALL : _____;
        icode == IJXX && Cnd: _____;
        icode == IRET : _____;
        1: _____;
    ];
        \end{minted}
        \sol 此题顺着上一题的填写就很容易写出了。
        \begin{minted}[frame=single, fontsize=\small]{text}
    int new_pc = [
        icode == ICALL : valC;
        icode == IJXX && Cnd: valC;
        icode == IRET : valM;
        1: valP;
    ];
        \end{minted}
        \pro 判断下列说法的正确性：
            \qn 流水线的深度越深，总吞吐率越大，因此流水线应当越深越好。
            \qn 流水线的吞吐率取决于最慢的流水级，因此流水线的划分应当尽量均匀。\qn 假设寄存器延迟为 \SI{20}{ps}，那么总吞吐率不可能达到或超过 \SI{50}{GIPS}。
            \qn 数据冒险总是可以只通过转发来解决。
            \qn 数据冒险总是可以只通过暂停流水线来解决。
        \sol 分别是错误（深度会影响冒险问题和电路复杂性，以及吞吐量有上界）、正确、正确、错误（load-use 冒险）、正确。
        \pro 一条三级流水线，包括延迟为 \SI{50}{ps}、\SI{100}{ps}、\SI{100}{ps} 的三个流水级，每个寄存器的延迟为 \SI{10}{ps}。那么这条流水线的总延迟是 \rule{2.5cm}{0.25mm} ps，吞吐率是 \rule{2.5cm}{0.25mm} GIPS。
        \sol 总延迟为 $100+100+100+3 \times 10 = \SI{330}{ps}$，吞吐率为 $1/330 = \SI{9.09}{GIPS}$。
        \pro 在下图中，A\textasciitilde H 为 8 个基本逻辑单元，图中标出了每个单元的延迟，以及用箭头标出了单元之间的数据依赖关系。寄存器的延迟均为 \SI{10}{ps}。
        \begin{figure}[H]
            \centering
            \begin{tikzpicture}
                \draw[draw=black] (0,0) rectangle ++(1,1.5);
                \node[align=center] at (0.5,0.75) {A \\ \SI{40}{ps}};
                \draw[draw=black] (2,0) rectangle ++(1,1.5);
                \node[align=center] at (2.5,0.75) {B \\ \SI{30}{ps}};
                \draw[draw=black] (4,0) rectangle ++(1,1.5);
                \node[align=center] at (4.5,0.75) {C \\ \SI{40}{ps}};
                \draw[draw=black] (6,0) rectangle ++(1,1.5);
                \node[align=center] at (6.5,0.75) {D \\ \SI{20}{ps}};
                \draw[-{Stealth[angle'=45]}] (1,0.75) -- (2,0.75);
                \draw[-{Stealth[angle'=45]}] (3,0.75) -- (4,0.75);
                \draw[-{Stealth[angle'=45]}] (5,0.75) -- (6,0.75);
                \draw[-{Stealth[angle'=45]}] (7,0.75) -- (8,0.75);

                \begin{scope}[shift={(0,-2.5)}]
                    \draw[draw=black] (0,0) rectangle ++(1,1.5);
                    \node[align=center] at (0.5,0.75) {E \\ \SI{20}{ps}};
                    \draw[draw=black] (2,0) rectangle ++(1,1.5);
                    \node[align=center] at (2.5,0.75) {F \\ \SI{60}{ps}};
                    \draw[draw=black] (4,0) rectangle ++(1,1.5);
                    \node[align=center] at (4.5,0.75) {G \\ \SI{40}{ps}};
                    \draw[draw=black] (6,0) rectangle ++(1,1.5);
                    \node[align=center] at (6.5,0.75) {H \\ \SI{30}{ps}};
                    \draw[-{Stealth[angle'=45]}] (1,0.75) -- (2,0.75);
                    \draw[-{Stealth[angle'=45]}] (3,0.75) -- (4,0.75);
                    \draw[-{Stealth[angle'=45]}] (5,0.75) -- (6,0.75);
                    \draw[-{Stealth[angle'=45]}] (7,0.75) -- (8,0.75);
                \end{scope}

                \draw[draw=black] (8,-2.5) rectangle ++(0.5,4);
                \node[align=center] at (8.25,-0.5) {R \\ E \\ G};
                \draw[-{Stealth[angle'=45]}] (1,0) -- (2,-1);
                \draw[-{Stealth[angle'=45]}] (5,-1) -- (6,0);
            \end{tikzpicture}
        \end{figure}
        \qn 计算目前的电路的总延迟。
        \qn 通过插入寄存器，可以对这个电路进行流水化改造。现在想将其改造为两级流水线， 为了达到尽可能高的吞吐率，问寄存器应插在何处？获得的吞吐率是多少？
        \qn 现在想将其改造为三级流水线，问最优改造所获得的吞吐率是多少？
        \sol 这里有多条数据通路，其总延迟由最高延迟的通路 $\mathrm A \to \mathrm F \to \mathrm G \to \mathrm H$ 决定，延迟是 $40+60+40+30+10 = \SI{180}{ps}$。插入寄存器使得其成为二级流水线，要求每一个通路都成为二级流水线，相当于多条单通路流水线求 $\max \min$。这里，通路 A\textasciitilde D 应当在 BC 之间插入寄存器；通路 E\textasciitilde H 应当在 FG 之间插入寄存器，此时各条数据通路都已经是二级的，最高单级延迟为 $40+60+10 = \SI{110}{ps}$，吞吐量 \SI{9.09}{GIPS}。对于三级流水线也是同样的道理，先考虑最慢通路 AFGH，它的两个寄存器应该放到 AF 和 FG 之间，此时 EFGH 的寄存器只能插在 EF 间（否则 AFGH 成为四级）。ABCD 通路的寄存器应插在 AB、BC 之间。综上，插入位置是 AB、AF、EF、BC、FG 之间，最高单级延迟 \SI{80}{ps}，吞吐量 \SI{12.5}{GIPS}。
    \end{problems}

\chapter{体系结构初步{---}往年考题}
    \begin{problems}
        \proy{2018} Y86 指令 \verb|popl rA| 的 SEQ 实现中，在取指阶段，\verb|valP <- |（\quad），在执行阶段，\verb|valE <- |（\quad）。
        \begin{choices}
            \item \verb|PC+4|、\verb|valA+4|
            \item \verb|PC+4|、\verb|valA+(-4)|
            \item \verb|PC+2|、\verb|valB+4|
            \item \verb|PC+2|、\verb|valB+(-4)|
        \end{choices}
        \proy{2017} 在 Y86 的 SEQ 实现中，对仅考虑 \verb|IRMMOVQ|、\verb|ICALL|、\verb|IPOPQ|、\verb|IRET| 指令， 对 \verb|mem_addr| 的 HCL 描述正确的是：
        \begin{minted}[frame=single, fontsize=\small]{text}
    word mem_addr = [
        icode in { (1), (2) } : valE;
        icode in { (3), (4) } : valA;
    ];
        \end{minted}
        \begin{choices}
            \item \verb|(1) IRMMOVQ  (2) IPOPQ  (3) IRET     (4) ICALL|
            \item \verb|(1) IRMMOVQ  (2) IRET   (3) IPOPQ    (4) ICALL|
            \item \verb|(1) ICALL    (2) IPOPQ  (3) IRMMOVQ  (4) IRET|
            \item \verb|(1) IRMMOVQ  (2) ICALL  (3) IPOPQ    (4) IRET|
        \end{choices}
        \proy{2017} 关于流水线技术的描述，错误的是: 
        \begin{choices}
            \item 流水线技术能够提高执行指令的吞吐率，但也同时增加单条指令的执行时间。
            \item 增加流水线级数，不一定能获得总体性能的提升。
            \item 指令间数据相关引发的数据冒险，不一定可以通过暂停流水线来解决。
            \item 流水级划分应尽量均衡，吞吐率会受到最慢的流水级影响，均衡的流水线能提高吞吐量。
        \end{choices}
        \proy{2017} 分析 64 位的 Y86 ISA 中新加入的条件内存传送指令：\verb|crmmovqXX| 和 \verb|cmrmovqXX|。\verb|crmmovqXX| 和 \verb|cmrmovqXX| 指令在条件码满足所需要的约束时，分别执行和 \verb|rmmovq| 以及 \verb|mrmovq| 同样的语义。其格式如下：
        \begin{table}[H]
            \tt
            \centering
            \begin{tabular}{c|c|c|c|c|c|}
                \cline{2-6}
                rmmovq & 4 & 0 & rA & rB & D (8 bytes) \\ \cline{2-6} 
                crmmovqXX & 4 & fn & rA & rB & D (8 bytes) \\ \cline{2-6} 
                mrmovq & 5 & 0 & rA & rB & D (8 bytes) \\ \cline{2-6} 
                cmrmovqXX & 5 & fn & rA & rB & D (8 bytes) \\ \cline{2-6} 
            \end{tabular}
        \end{table}
        \qn 请按下表补全每个阶段的操作。需说明的信号可能会包括：\verb|icode|、\verb|ifun|、\verb|rA|、\verb|rB|、\verb|valA|、\verb|valB|、\verb|valC|、\verb|valE|、\verb|valP|、\verb|Cnd|；寄存器堆 \verb|R[]|、存储器 \verb|M[]|、程序计数器 \verb|PC|、条件码 \verb|CC|。其中对存储器的引用必须标明字节数。
        \begin{table}[H]
            \centering
            \begin{tabular}{|c|c|c|}
                \hline
                阶段 & {\qquad} \verb|rmmovq rA, D(rB)| {\qquad} & {\qquad} \verb|cmrmovqXX D(rB), rA| {\qquad} \\ \hline
                取指 & \multicolumn{2}{c|}{\rule{0pt}{10ex}} \\ \hline
                译码 & \multicolumn{2}{c|}{\texttt{valA <- R[rA], valB <- R[rB]}} \\ \hline
                执行 & \rule{0pt}{10ex} & \rule{0pt}{10ex} \\ \hline
                访存 & \rule{0pt}{10ex} & \rule{0pt}{10ex} \\ \hline
                写回 & N/A & \rule{0pt}{10ex} \\ \hline
                更新 \verb|PC| & \multicolumn{2}{c|}{\texttt{PC <- valP}} \\ \hline
            \end{tabular}
        \end{table}
        \sol 取指部分的填法是很基本的，需要获得 \verb|rA, rB, D|，所以我们写 \verb|icode:ifun <- M1[PC]|、 \verb|rA:rB <- M1[PC+1]|、\verb|valC <- M8[PC+2]| 和 \verb|vavlP <- PC+10|。译码部分也很简单，为 \verb|valA <- R[rA]|、\verb|valB <- R[rB]|。

        在执行部分，\verb|rmmovq| 都要计算地址偏移，因此 \verb|valE <- valB+valC|，但是 \verb|cmrmovqXX| 还需要处理条件码，所以还要填 \verb|Cnd <- Cond(CC, ifun)|。最后 \verb|rmmovq| 还需要访存，即 \verb|M8[valE] <- valA|；而 \verb|cmrmovqXX| 除了访存 \verb|valM <- M8[valE]| 之外，还需要条件写回 \verb|if (Cnd) R[rA] <- valM|。
        \proy{2016} 下面对指令系统的描述中，错误的是：
        \begin{choices}
            \item CISC 指令系统中的指令数目较多，有些指令的执行周期很长；而 RISC 指令系统中通常指令数目较少，指令的执行周期都较短。
            \item CISC 指令系统中的指令编码长度不固定；RISC 指令系统中的指令编码长度固定，这样使得 CISC 机器可以获得了更短的代码长度。
            \item CISC指令系统支持多种寻址方式，RISC指令系统支持的寻址方式较少。
            \item CISC 机器中的寄存器数目较少，函数参数必须通过栈来进行传递；RISC 机器中的寄存器数目较多，只需要通过寄存器来传递参数，避免了不必要的存储访问。
        \end{choices}
        \proy{2016} 下面对流水线技术的描述，正确的是：
        \begin{choices}
            \item 流水线技术不仅能够提高执行指令的吞吐率，还能减少单条指令的执行时间。
            \item 不断加深流水线级数，总能获得性能上的提升。
            \item 流水级划分应尽量均衡，吞吐率会受到最慢的流水级影响。
            \item 指令间的数据相关可能会引发流水线停顿，但总是可以通过调度指令来解决。
        \end{choices}
        \proy{2016} 若处理器实现了三级流水线，每一级流水线实际需要的运行时间分别为 \SI{1}{ns}、\SI{2}{ns} 和 \SI{3}{ns}，则此处理器不停顿地执行完毕 10 条指令需要的时间为：
        \begin{choices}
            \item \SI{21}{ns}
            \item \SI{12}{ns}
            \item \SI{24}{ns}
            \item \SI{36}{ns}
        \end{choices}
        \sol 首先，每级流水线都是 \SI{3}{ns}，然后画出时序图就可以了。一般地，需要时间是一级流水线的总时间，补上一头一尾。这里需要的时间为 $10 \times 3+3+3 = \SI{36}{ns}$。
        \proy{2016} 请分析 32 位的 Y86 ISA 中新加入的一组条件返回指令 \verb|cretXX|，其格式如下。
        \begin{table}[H]
            \tt
            \centering
            \begin{tabular}{c|c|c|}
                \cline{2-3}
                cretXX & 9 & fun \\ \cline{2-3} 
            \end{tabular}
        \end{table}
        类似 \verb|cmovXX|，该组指令只有当条件码 \verb|Cnd| 满足时，才执行函数返回；如果条件不满足，则顺序执行。

        若在教材所描述的 SEQ 处理器上执行这条指令，请按下表补全每个阶段的操作。需说明的信号可能会包括：\verb|icode|、\verb|ifun|、\verb|rA|、\verb|rB|、\verb|valA|、\verb|valB|、\verb|valC|、\verb|valE|、\verb|valP|、\verb|Cnd|；寄存器堆 \verb|R[]|、存储器 \verb|M[]|、程序计数器 \verb|PC|、条件码 \verb|CC|。其中对存储器的引用必须标明字节数。如果在某一阶段没有任何操作，请填写 \texttt{none} 指明。
        \begin{table}[H]
            \centering
            \begin{tabular}{|c|c|}
                \hline
                阶段 & {\qquad \qquad \qquad \qquad} \verb|cretXX offset| {\qquad \qquad \qquad \qquad} \\ \hline
                取指 & \rule{0pt}{10ex} \\ \hline
                译码 & \rule{0pt}{10ex} \\ \hline
                执行 & \rule{0pt}{10ex} \\ \hline
                访存 & \rule{0pt}{10ex} \\ \hline
                写回 & \rule{0pt}{10ex} \\ \hline
                更新 \verb|PC| & \rule{0pt}{10ex} \\ \hline
            \end{tabular}
        \end{table}
        \sol 本质还是 \verb|ret|，只是需要稍微添加一下条件判断。取指是 \verb|icode:ifun <- M1[PC]|、\verb|valP <- PC+1|。译码需要注意，在 \verb|ret| 中两个值都会被取用，故填写 \verb|valB <- R[%rsp]|、\verb|valA <- R[%rsp]|。
        
        执行阶段是要恢复栈空间，因此为 \verb|valE <- valB+4| 和 \verb|Cnd <- Cond(CC, ifun)|。这里用 \verb|valA| 也可以。访存阶段需要读出返回地址，请注意，这里必须用 \verb|valA|（参看SEQ 实现图中的数据通路），填写 \verb|valM <- M4[valA]|。随后的写回和更新程序计数器的工作主要是关心条件（是否真的要返回了），所以分别是 \verb|if (Cnd) R[%rsp] <- valE| 和 \verb|PC <- Cnd?valM:valP|。
        \proy{2015} 下面有关指令系统设计的描述正确的是：
        \begin{choices}
            \item 采用 CISC 指令比 RISC 指令代码更长。
            \item 采用 CISC 指令比 RISC 指令运行时间更短。
            \item 采用 CISC 指令比 RISC 指令译码电路更加复杂。
            \item 采用 CISC 指令比 RISC 指令的流水线吞吐更高。
        \end{choices}
        \proy{2015} 一个功能模块包含组合逻辑和寄存器，组合逻辑单元的总延迟是 \SI{100}{ps}，单个寄存器的延时是 20ps，该功能模块执行一次并保存执行结果，理论上能达到的最短延时和最大吞吐分别是多少？
        \begin{choices}
            \item \SI{20}{ns}、\SI{50}{GIPS}
            \item \SI{120}{ns}、\SI{50}{GIPS}
            \item \SI{120}{ns}、\SI{10}{GIPS}
            \item \SI{20}{ps}、\SI{10}{GIPS}
        \end{choices}
        \proy{2015} 关于流水线技术的描述，错误的是：
        \begin{choices}
            \item 流水线技术能够提高执行指令的吞吐率，但也同时增加单条指令的执行时间。
            \item 减少流水线的级数，能够减少数据冒险发生的几率。
            \item 指令间数据相关引发的数据冒险，都可以通过数据转发来解决。
            \item 现代处理器支持一个时钟内取指、执行多条指令，会增加控制冒险的开销。
        \end{choices}
        \proy{2015} 请分析 Y86 ISA 中新加入的一条指令：\verb|NewJE|，其格式如下。
        \begin{table}[H]
            \centering
            \begin{tabular}{c|c|c|c|c|c|}
                \cline{2-6}
                NewJE & C & 0 & rA & rB & dst \\ \cline{2-6} 
            \end{tabular}
        \end{table}
        其功能为：如果 \verb|R[rA] == R[rB]|，则跳转到 \verb|dst| 继续执行，否则顺序执行。

        若在教材所描述的 SEQ 处理器上执行这条指令，请按下表补全每个阶段的操作。需说明的信号可能会包括：\verb|icode|、\verb|ifun|、\verb|rA|、\verb|rB|、\verb|valA|、\verb|valB|、\verb|valC|、\verb|valE|、\verb|valP|、\verb|Cnd|；寄存器堆 \verb|R[]|、存储器 \verb|M[]|、程序计数器 \verb|PC|、条件码 \verb|CC|。其中对存储器的引用必须标明字节数。如果在某一阶段没有任何操作，请填写 \texttt{none} 指明。
        \begin{table}[H]
            \centering
            \begin{tabular}{|c|c|}
                \hline
                阶段 & {\qquad \qquad \qquad \qquad} \verb|cretXX offset| {\qquad \qquad \qquad \qquad} \\ \hline
                取指 & \rule{0pt}{10ex} \\ \hline
                译码 & \rule{0pt}{10ex} \\ \hline
                执行 & \rule{0pt}{10ex} \\ \hline
                访存 & \rule{0pt}{10ex} \\ \hline
                写回 & \rule{0pt}{10ex} \\ \hline
                更新 \verb|PC| & \rule{0pt}{10ex} \\ \hline
            \end{tabular}
        \end{table}
        \proy{2014} 若处理器实现了三级流水线，每一级流水线实际需要的运行时间分别为 \SI{2}{ns}、\SI{2}{ns} 和 \SI{1}{ns}，则此处理器不停顿地执行完毕 10 条指令需要的时间为：
        \begin{choices}
            \item \SI{21}{ns}
            \item \SI{22}{ns}
            \item \SI{23}{ns}
            \item \SI{24}{ns}
        \end{choices}
        \proy{2014} 关于 RISC 和 CISC 的描述，正确的是：
        \begin{choices}
            \item CISC 指令系统的指令编码可以很短，例如最短的指令可能只有一个字节，因此 CISC 的取指部件设计会比 RISC 更为简单。
            \item CISC 指令系统中的指令数目较多，因此程序代码通常会比较长；而 RISC 指令系统中通常指令数目较少，因此程序代码通常会比较短。
            \item CISC 指令系统支持的寻址方式较多，RISC 指令系统支持的寻址方式较少，因此用 CISC 在程序中实现访存的功能更容易。
            \item CISC 机器中的寄存器数目较少，函数参数必须通过栈来进行传递；RISC 机器中的寄存器数目较多，只需要通过寄存器来传递参数。
        \end{choices}
        \proy{2014} 关于流水线技术的描述，正确的是：
        \begin{choices}
            \item 指令间数据相关引发的数据冒险，一定可以通过暂停流水线来解决。
            \item 流水线技术不仅能够提高执行指令的吞吐率，还能减少单条指令的执行时间。
            \item C. 增加流水线的级数，一定能获得性能上的提升。
            \item 流水级划分应尽量均衡，不均衡的流水线会增加控制冒险。
        \end{choices}
        \proy{2014} 请分析 Y86 ISA 中新加入的一条指令：\verb|caddXX| 条件加法。其功能可以参考 \verb|add| 和 \verb|cmovXX| 两条指令。
        \begin{table}[H]
            \centering
            \begin{tabular}{c|c|c|c|c|}
                \cline{2-5}
                caddXX & C & fn & rA & rB \\ \cline{2-5} 
            \end{tabular}
        \end{table}

        若在教材所描述的 SEQ 处理器上执行这条指令，请按下表补全每个阶段的操作。需说明的信号可能会包括：\verb|icode|、\verb|ifun|、\verb|rA|、\verb|rB|、\verb|valA|、\verb|valB|、\verb|valC|、\verb|valE|、\verb|valP|、\verb|Cnd|；寄存器堆 \verb|R[]|、存储器 \verb|M[]|、程序计数器 \verb|PC|、条件码 \verb|CC|。其中对存储器的引用必须标明字节数。如果在某一阶段没有任何操作，请填写 \texttt{none} 指明。
        \begin{table}[H]
            \centering
            \begin{tabular}{|c|c|}
                \hline
                阶段 & {\qquad \qquad \qquad \qquad} \verb|cretXX offset| {\qquad \qquad \qquad \qquad} \\ \hline
                取指 & \rule{0pt}{10ex} \\ \hline
                译码 & \rule{0pt}{10ex} \\ \hline
                执行 & \rule{0pt}{10ex} \\ \hline
                访存 & \rule{0pt}{10ex} \\ \hline
                写回 & \rule{0pt}{10ex} \\ \hline
                更新 \verb|PC| & \rule{0pt}{10ex} \\ \hline
            \end{tabular}
        \end{table}
    \end{problems}
	\chapter{存储器层次结构}
    \begin{summary}
        \begin{compactitem}
            \item 了解 SRAM、DRAM 等易失性存储器及其变种（如 DDR SDARM 等）的功能和历史；了解 PROM、EPROM 等非易失性存储器及其变种的功能和历史。
            \item 知道传统旋转磁盘的结构相关概念和工作原理，会计算其存储容量和使用磁盘访问数据的延迟。了解固态硬盘的工作原理。
            \item 了解计算机体系结构中总线的作用以及含总线的计算机基本结构。
            \item 熟悉时间局部性和空间局部性的概念，会阅读程序判断局部性的概念。
        \end{compactitem}
    \end{summary}

    \begin{problems}
        \pro 对于下列描述，是 SRAM 更符合还是 DRAM 更符合，还是均符合？
            \qn 访问速度更快
            \qn 每比特需要的晶体管数目少
            \qn 单位容量造价更便宜
            \qn 常用作主存
            \qn 需要定期刷新
            \qn 断电后失去存储的信息
            \qn 支持随机访问
            \qn 一种变种是 SDRAM
        \sol SRAM 较快，复杂且贵，而 DRAM 相对反之。所以 1 属于 SRAM 而 2、3、4、5 都是 DRAM。6、7 都是 RAM 的属性，而 SDRAM 是一种 DRAM 的变种（SDRAM 是同步动态随机存取存储器的缩写）。
        \pro 下列存储器中，属于易失性存储介质的有：
            \begin{choices}
                \item DRAM
                \item SRAM
                \item ROM
                \item 软盘
                \item SSD
                \item U 盘
            \end{choices}
        \sol AB，属于基本常识题。
        \pro 已知一个双面磁盘有 2 个盘片、10000 个柱面，每条磁道有 400 个扇区，每个扇区容量为 512 字节，则它的存储容量是 \rule{2.5cm}{0.25mm} GB。
        \sol 因为是双面磁盘，所以计算 $2 \times 2 \times 10000 \times 400 \times 512 = \SI{8192000000}{Byte} = \SI{8.192}{GB}$，注意一般较慢的存储介质（例如磁盘容量、网速等）都是用十进制的数前缀。
        \pro 已知一个磁盘的平均寻道时间为 \SI{6}{ms}，旋转速度为 \SI{7500}{RPM}，那么它的平均访问时间大约为 \rule{2.5cm}{0.25mm} ms。
        \sol 平均地看，需要旋转半周。又磁盘主要的延迟来源于寻道和旋转，所以结果为 $6 + 0.5 \times (60/7500 \times 1000) = \SI{10}{ms}$。
        \pro 已知一个磁盘每条磁道平均有 400 个扇区，旋转速度为 \SI{6000}{RPM}，那么它的平均传送时间大约为 \rule{2.5cm}{0.25mm} ms。
        \sol 答案为 $60/6000 \text{ （转一圈的时间）} \times 1/400 \text{ （转过一个磁道的时间）} \times 1000 = \SI{0.025}{ms}$。
        \pro 考虑如下程序：
        \begin{minted}[frame=single, fontsize=\small]{c}
    for (int i = 0; i < n; i++) {
        B[i] = 0;
        for (int j = 0; j < m; j++)
            B[i] += A[i][j];
    }
        \end{minted}
        判断下列说法的正确性。
            \qn 对于数组 \verb|A| 的访问体现了时间局部性。
            \qn 对于数组 \verb|A| 的访问体现了空间局部性。
            \qn 对于数组 \verb|B| 的访问体现了时间局部性。
            \qn 对于数组 \verb|B| 的访问体现了空间局部性。
        \sol 数组 \verb|A, B| 都存在逐元素访问的过程，因此都体现了空间局部性。但 \verb|B| 存在同一个元素反复访问而 \verb|A| 没有。因此 1 错误而 2、3、4 均正确。顺便一提：任何一个有意义的程序的指令都体现时间局部性（程序计数器）和空间局部性（顺序执行）。
        \pro 回答下列有关缓存结构的问题：
            \qn 一个容量为 \SI{8}{K} 的直接映射高速缓存，每一行的容量为 \SI{32}{B}，那么它有 \rule{2.5cm}{0.25mm} 组，每组有 \rule{2.5cm}{0.25mm} 行。
            \qn 一个容量为 \SI{8}{K} 的全相联高速缓存，每一行的容量为 \SI{32}{B}，那么它有 \rule{2.5cm}{0.25mm} 组，每组有 \rule{2.5cm}{0.25mm} 行。
            \qn 一个容量为 \SI{16}{K} 的 4 路组相联告诉缓存，每一行的容量为 \SI{64}{B}，那么一个 16 位地址 \verb|0xCAFE| 应映射在第 \rule{2.5cm}{0.25mm} 组内。
        \sol 总容量等于组数乘以路数（行数）乘以块大小，其中直接映射高速缓存路数为 1，而全相联高速缓存的组数为 1。因此第一问为 $2^3 \times 2^10/2^5=256$ 组，1 行；第二问为 1 组，$2^3 \times 2^10/32=256$ 行。对于第三问，组数为 $2^4 \times 2^10/(2^6 \times 4)=64$ 组，所以组索引占据 6 位，偏移量占据 6。所给出的地址为 \verb|1100 1010 1111 1110|，即组索引是 \verb|101011|，即 43。
        \pro 判断下列说法的正确性。
            \qn 保持块大小与路数不变，增大组数，命中率一定不会降低。
            \qn 保持总容量与块大小不变，增大路数，命中率一定不会降低。
            \qn 保持总容量与路数不变，增大块大小，命中率一定不会降低。
            \qn 使用随机替换代替 LRU，期望命中率可能会提高。
        \sol 1 正确，相当于在 cache 中直接增加一组空间。2、3 都错误，由于总容量不变，所以组数都会降低，命中率都可能降低。实际上实验表明，增大路数一般会导致命中率降低，而增大块大小对于不同的访存步长可能造成命中率升高或降低。4 正确，对于某些特定的程序，LRU 的期望性能可能不如随机替换（采看 2018 年期中第 13 题的构造），实验表明在较小的 cache 中随机策略更优。
        \pro 考虑如下程序：
        \begin{minted}[frame=single, fontsize=\small]{c}
    int A[MAXN];
    for (int i = 0; i < 25; i++) {
        int x = A[i];
        int y = A[i + 1];
        int z = A[i + 2];
        A[i + 3] = x + y + z;
    }
        \end{minted}
        假设编译成汇编语言的时候没有任何优化，变量 \verb|x|、\verb|y|、\verb|z| 均放在寄存器中，运行之前 cache 所有行都是无效的。\verb|A| 的起始地址为 \verb|0|，且 \verb|MAXN| 充分大，保证没有越界问题。
        \qn 假设 cache 的容量为 8 字节，每一行的容量为 4 字节，替换策略为 LRU，组策略为直接映射高速缓存。在这个 cache 上运行上述代码，得到的 cache 命中率是 \rule{2.5cm}{0.25mm}\%。
        \qn 假设 cache 的容量为 8 字节，每一行的容量为 4 字节，替换策略为 LRU，组策略为全相联高速缓存。在这个 cache 上运行上述代码，得到的 cache 命中率是 \rule{2.5cm}{0.25mm}\%。
        \qn 假设 cache 的容量为 32 字节，每一行的容量为 8 字节，替换策略为 LRU，组策略为 2 路组相联。画出程序运行结束时 cache 的情况（用 \verb|M[0-7]| 表示第 \verb|0| 到第 \verb|7| 字节的地址）；cache 命中率是 \rule{2.5cm}{0.25mm}\%。
        \qn 假设 cache 的每一行的容量为 4 字节，运行该程序，得到的 cache 命中率的可能最大值为 \rule{2.5cm}{0.25mm}\%。
        \sol 处理 cache 命中率的题，首先确定 tag, index 和 offset 的分布，然后先尝试通过组索引简化分析，\emph{获得各地址放置位置的规律}，最后再直接模拟。对于本题，首先知道每次访问 4 个整数，cache 的每行（路）都能容纳 1 个元素或者 2 个元素（视小问而不同）。一共访存 100 次，只需要计算命中的次数。我们下面用下标 0、1、2、3... 来代指数组中的元素。

        第一问。计算得一共 2 组，每组 1 行，所以最低三位中，高 1 位是组索引，剩下 2 位是偏移。因为是直接映射，且每行只能容纳一个整数，所以组索引其实也代表了元素在数组的下标，因此实际上两行分别存储下标为奇数和偶数的元素。第一次 0、1、2、3 均不命中，此时 2、3 留在缓存中。接着访问 1、2、3、4，但 3 会被 1 替换，只有 2 能命中，结束时 3、4 留在缓存中。再次访问 2、3、4、5，但 4 会被 2 替换，只有 3 能命中。由此发现后 24 组访问各恰有一次命中，命中率为 24\%。

        第二问。计算得只有 1 组，组内有 2 行，因此只有最低 2 位代表组内偏移，每行只能容纳一个整数，因此是依次填满并 LRU 替换。第一次 0、1、2、3 均不命中，此时 2、3 留在缓存中。接着访问 1、2、3、4 时，首先 1 不命中，替换 2，然后 2 不命中，替换 3，最后 4 不命中，替换 2，以此类推。由此发现这个访存序列发生了抖动，每次都不命中，命中率为 0\%。

        第三问需要一些细致的模拟。计算得有 2 组，每组 2 行，每行都能容纳两个连续的整数。在这里，组索引表示的是下标模 4 的结果，也就是说 0、1、4、5... 会进入一组，而 2、3、6、7... 进入另一组。我们预期访存序列模 4 会有某种规律。模拟过程如下（你可能没必要模拟这么长）：
        \begin{compactenum}[(i)]
            \item\ 0、1、2、3，其中 0、2 不命中，结束时第一组有 0、1，第二组有 2、3。
            \item\ 1、2、3、4，其中 4 不命中，结束时第一组有 0、1、4、5，第二组有 2、3。
            \item\ 2、3、4、5，全部命中。
            \item\ 3、4、5、6，只有 6 不命中，结束时第一组有 0、1、4、5，第二组有 2、3、6、7。
            \item\ 4、5、6、7，全部命中。
            \item\ 5、6、7、8，只有 8 不命中，此时开始替换，第一组有 8、9、4、5，第二组有 2、3、6、7。
            \item\ 6、7、8、9，全部命中。
            \item\ 7、8、9、10，只有 10 不命中，再次替换，第一组有 8、9、4、5，第二组有 10、11、6、7。
            \item\ 8、9、10、11，全部命中。
            \item\ 9、10、11、12，只有 12 不命中，替换得到第一组有 8、9、12、13，第二组有 10、11、6、7。
        \end{compactenum}
        这样就大致找到规律了，不命中只可能是在 $0, 2, \dotsc, 24, 26$ 时发生，这里一共有 14 次，因此命中率为 86\%。最后一次访存是 24、25、26、27，是全部命中的，根据(iv)、(vi)、(ix)的规律，其中剩下 cache 的部分应该对应前四个元素，即 20 \textasciitilde23 都在缓存中。综上所述，地址 \verb|80| 起的 8 个字节都在缓存中，所以缓存可以画为
        \begin{table}[H]
            \tt
            \centering
            \begin{tabular}{|c|c|c|c|c|}
                \hline
                & 有效位 & 内容 & 有效位 & 内容 \\ \hline
                组 0 & 1 & M[96-103] & 1 & M[80-87] \\ \hline
                组 1 & 1 & M[88-95] & 1 & M[104-111] \\ \hline
            \end{tabular}
        \end{table}

        第四问考察冷不命中。当缓存充分大时，至少第一次访问该元素时是要发生不命中的，而这里一共有 28 个元素要访问，所以最高命中率是 72\%。
    \end{problems}

\chapter{存储器层次结构{---}往年考题}
    \begin{problems}
        \proy{2018} 以下关于存储的叙述中，正确的是：
        \begin{choices}
            \item 由于基于 SRAM 的内存性能与 CPU 的性能有很大差距，因此现代计算机使用更快的基于 DRAM 的高速缓存，试图弥补 CPU 和内存间性能的差距。
            \item SSD 相对于旋转磁盘而言具有更好的读性能，但是 SSD 写的速度通常比读的速度慢得多，而且 SSD 比旋转磁盘单位容量的价格更贵，此外 SSD 底层基于 EEPROM 的闪存会磨损。
            \item 一个有 2 个盘片，10000 个柱面，每条磁道平均有 400 个扇区，每个扇区有 512 个字节的双面磁盘的容量为 \SI{8}{GB}。
            \item 访问一个磁盘扇区的平均时间主要取决于寻道时间和旋转延迟，因此一个旋转速率为 \SI{6000}{RPM}、平均寻道时间为 \SI{9}{ms} 的磁盘的平均访问时间大约为 \SI{19}{ms}。
        \end{choices}
        \proy{2017} 以下计算机部件中，通常不用于存储器层次结构（Memory Hierarchy）的是：
        \begin{choices}
            \item 高速缓存
            \item 内存
            \item 硬盘
            \item 优盘（U 盘）
        \end{choices}
        \proy{2017} 关于局部性（locality）的描述，正确的是：
        \begin{choices}
            \item 数据的时间局部性或数据空间局部性，在任何有意义的程序中都能体现。
            \item 指令的时间局部性或数据空间局部性，在任何有意义的程序中都能体现。
            \item 数据的时间局部性，在任何循环操作中都能体现。
            \item 数据的空间局部性，在任何数组操作中都能体现。
        \end{choices}
        \proy{2016} 以下关于存储结构的讨论，哪个是正确的？
        \begin{choices}
            \item 增加额外一级存储，数据存取的延时一定不会下降。
            \item 增加存储的容量，数据存取的延时一定不会下降。
            \item 增加额外一级存储，数据存取的延时一定不会增加。
            \item 以上选项都不正确。
        \end{choices}
        \proy{2016} 关于局部性（locality）的描述，不正确的是：
        \begin{choices}
            \item 循环通常具有很好的时间局部性。
            \item 循环通常具有很好的空间局部性。
            \item 数组通常具有很好的时间局部性。
            \item 数组通常具有很好的空间局部性。
        \end{choices}
        \proy{2015} 下面关于存储器的说法，错误的是：
        \begin{choices}
            \item SDRAM 的速度比 FPM DRAM 快。
            \item SDRAM 的 RAS 和 CAS 请求共享相同的地址引脚。
            \item 磁盘的寻道时间和旋转延迟大致在一个数量级。
            \item 固态硬盘的随机读写性能基本相当。
        \end{choices}
        \proy{2015} 某磁盘的旋转速率为 \SI{7200}{RPM}，每条磁道平均有 400 扇区，则一个扇区的平均传送时间为 \rule{2.5cm}{0.25mm} ms。
        \proy{2018} 一缓存共包含 4 个缓存块，每个块 32 字节，采用 LRU 替换算法。设缓存初始为空，对其进行下列 \verb|char| 型内存地址序列访存：\verb|0x5a7, 0x5b7, 0x6a6, 0x5b8, 0x7a5, 0x5b9|。对于直接映射缓存和 2 路组相联缓存，分别能产生几次命中？
        \begin{choices}
            \item 1, 3
            \item 1, 4
            \item 2, 3
            \item 2, 4
        \end{choices}
        \proy{2018} 设一种全相联缓存共包含 4 个缓存块，如果循环地顺序访问 5 个不同的内存块（大小和一个缓存块一样），下列哪种替换算法会产生最多的命中？（考虑渐近值即可）
        \begin{choices}
            \item 最近最少使用替换策略 LRU
            \item 先入先出替换策略 FIFO
            \item 随机替换策略 Random
            \item 后入先出替换策略 LIFO
        \end{choices}
        \sol CD。对于全相联缓存，操作主要和替换策略有关。LRU 和 FIFO 的行为是一样的，并且发生抖动（第 $n$ 次访问会替换第 $n+1$ 次访问需要的块），因此不会发生任何命中。随机替换策略中，每次新来一个块，随机换出的块会导致之后的四次访存有一次 miss，且是等概率地有一次。因而平均 2.5 次访存会发生一次 miss，所以命中率是 60\%。LIFO 策略中，前 3 个内存块的访问一直会命中，而最后两个块交替发生 miss，所以命中率也是 60\%。
        \proy{2018} 某计算机地址空间 12 位，L1 cache 大小为 256 字节，组数 $S=4$，路数 $E=2$，现在在地址 \verb|0x0| 处开始有一个 $N$ 行 $M$ 列的 \verb|int| 类型的数组 \verb|int A[N][M]|，有如下 C 代码：
        \begin{minted}[frame=single, fontsize=\small]{c}
    int ans = 0;
    for(int j = 0; j < M; ++j)
        for(int i = 0; i < N; ++i) ans += A[i][j];
        \end{minted}
        不考虑并行、编译优化等等会影响访问数组元素顺序的因素，只使用 L1 cache，则如下 $(N, M)$ 对中会导致全部 cache miss 的有几个？$(64, 4), (32, 8), (16, 16), (2, 128)$
        \begin{choices}
            \item 4 个
            \item 3 个
            \item 2 个
            \item 1 个
        \end{choices}
        \sol C。计算得到每个缓存块可以存放 $256/(4SE)=8$ 个整数 $(N, M)=(64, 4)$ 时，访问数组第一行元素会将第一、二行元素都放入缓存，因此在访问数组第二行元素 \verb|A[1][0]| 时不会 miss；而 $(N, M)=(2, 128)$ 时，同理可得访问 \verb|A[0][1]| 时不会 miss。利用组索引简化分析，可以验证 $(N, M)=(32, 8), (16, 16)$ 时都会发生全部 miss。
        \proy{2018} Cache 为处理器提供了一个高性能的存储器层次框架。下面是一个 8 位存储器地址引用的列表（地址单位为字节，数字为 10 进制）：\texttt{3, 180, 43, 2, 191, 88, 190, 14, 181, 44}。
        \qn 考虑如下 cache：$S=2, E=2$，块大小为 2 字节；初始状态为空，替换策略为 LRU。经历上面的访问序列后，请在下表中填写缓存的状态。格式要求：tag 使用二进制；data 使用十进制，例如 \verb|M[6-7]| 表示地址 6 和 7 对应的数据。
        \begin{table}[H]
            \tt
            \centering
            \begin{tabular}{|c|c|c|c|c|c|c|}
                \hline
                & {\quad V \quad} & {\quad Tag \quad} & {\quad Data \quad} & {\quad V \quad} & {\quad Tag \quad} & {\quad Data \quad} \\ \hline
                SET 0 &  &  &  &  &  &  \\ \hline
                SET 1 &  &  &  &  &  &  \\ \hline
            \end{tabular}
        \end{table}
        \qn 现在有另外两种直接映射的 cache 设计方案 C1 和 C2，每种方案的 cache 总大小都为 8 字节，C1 块大小为 2 字节，C2 块大小为 4 字节。假设从内存加载一次数据到 cache 的时间为 25 个周期，访问一次 C1 的时间为 3 个周期，访问一次 C2 的时间为 5 个周期。针对(1)的地址访问序列，哪一种 cache 的设计更好？请分别给出两种 cache 访问第一问地址序列的总时间以及 miss rate。
        \qn 现在考虑另外一个计算机系统。在该系统中，存储器地址为 32 位，并采用容量 \SI{32}{KB}、块大小 8 字节的直接映射高速缓存，则此 cache 实际至少要占用 \rule{2.5cm}{0.25mm} 字节的空间。（$\text{datasize} + (\text{valid bit size} + \text{tag size}) \times \text{blocks}$）
        \proy{2017} 在高速缓存存储器中，关于全相联和直接映射结构，以下论述正确的是：
        \begin{choices}
            \item 如果配备同样容量、技术的高速缓存，配备全相联高速缓存的计算机总是比配备直接映射高速缓存的计算机性能低。
            \item 如果配备同样容量、技术的高速缓存，配备全相联高速缓存的计算机总是比配备直接映射高速缓存的计算机性能高。
            \item 如果配备同样容量、技术的高速缓存，当数据在缓存中时，配备全相联高速缓存的计算机总是比配备直接映射高速缓存的计算机读数据慢。
            \item 如果配备同样容量、技术的高速缓存，当数据在缓存中时，配备全相联高速缓存的计算机总是比配备直接映射高速缓存的计算机读数据快。
        \end{choices}
        \proy{2017} 现有一个能够存储 4 个 block 的 cache，每一个 cache block 的大小为 \SI{2}{Byte}（$B=2$）。内存空间的大小是 \SI{32}{Byte}，即其地址范围为 \texttt{0\textsubscript{10} (00000\textsubscript{2}) --- 31\textsubscript{10} (11111\textsubscript{2})}。此设定下一程序的访问内存地址序列如下所示，单位是 Byte，数字为十进制：
        \begin{center}
            \tt 0, 3, 4, 7, 16, 19, 21, 22, 8, 10, 13, 14, 24, 26, 29, 30.
        \end{center}
        \qn 若 cache 的结构如下图所示（$S=2, E=2$），初始状态为空，替换策略为 LRU。请在下图空白处填入上述数据访问后 cache 的状态。格式要求：tag 使用二进制；data 使用十进制，例如 \verb|M[6-7]| 表示地址 6 和 7 对应的数据。
        \begin{table}[H]
            \tt
            \centering
            \begin{tabular}{|c|c|c|c|c|c|c|}
                \hline
                & {\quad V \quad} & {\quad Tag \quad} & {\quad Data \quad} & {\quad V \quad} & {\quad Tag \quad} & {\quad Data \quad} \\ \hline
                SET 0 &  &  &  &  &  &  \\ \hline
                SET 1 &  &  &  &  &  &  \\ \hline
            \end{tabular}
        \end{table}
        上述数据访问一共产生了 \rule{2.5cm}{0.25mm} 次命中。
        \qn 在(1)的基础上增加一条数据预取规则：每当 cache 访问出现 miss 时，被访问地址及其后续的一个 cache block 都会被放入缓存，例如当 \verb|M[0-1]| 访问发生 miss，则把 \verb|M[0-1]| 和 \verb|M[2-3]| 都放入缓存中。此时这 16 次数据访问一共产生了 \rule{2.5cm}{0.25mm} 次命中。
        \qn 在(1)的基础上将每一个 cache block 的大小扩大为 \SI{4}{Byte}（即 $B=4$，cache 大小变为原来的 2 倍），此时这 16 次数据访问一共产生了 \rule{2.5cm}{0.25mm} 次命中。
        \qn 在(3)的基础上，考虑增加(2)中的数据预取规则，则这 16 次数据访问一共产生了 \rule{2.5cm}{0.25mm} 次命中。
        \proy{2016} 现有一个能够存储 4 个 block 的 cache，每一个 cache block 的大小为 \SI{2}{Byte}（$B=2$）。内存空间的大小是 \SI{32}{Byte}，即其地址范围为 \texttt{0\textsubscript{10} (00000\textsubscript{2}) --- 31\textsubscript{10} (11111\textsubscript{2})}。此设定下一程序的访问内存地址序列如下所示，单位是 Byte，数字为十进制：
        \begin{center}
            \tt 2, 23, 10, 9, 9, 11, 3.
        \end{center}
        \qn 若 cache 的结构如下图所示（$S=2, E=2$），初始状态为空，替换策略为 LRU。请在下图空白处填入上述数据访问后 cache 的状态。格式要求：tag 使用二进制；data 使用十进制，例如 \verb|M[6-7]| 表示地址 6 和 7 对应的数据。
        \begin{table}[H]
            \tt
            \centering
            \begin{tabular}{|c|c|c|c|c|c|c|}
                \hline
                & {\quad V \quad} & {\quad Tag \quad} & {\quad Data \quad} & {\quad V \quad} & {\quad Tag \quad} & {\quad Data \quad} \\ \hline
                SET 0 &  &  &  &  &  &  \\ \hline
                SET 1 &  &  &  &  &  &  \\ \hline
            \end{tabular}
        \end{table}
        上述数据访问一共产生了 \rule{2.5cm}{0.25mm} 次命中。
        \qn 若替换策略为 MRU，其余和(1)相同，请在下图空白处填入上述数据访问后 cache 的状态。
        \begin{table}[H]
            \tt
            \centering
            \begin{tabular}{|c|c|c|c|c|c|c|}
                \hline
                & {\quad V \quad} & {\quad Tag \quad} & {\quad Data \quad} & {\quad V \quad} & {\quad Tag \quad} & {\quad Data \quad} \\ \hline
                SET 0 &  &  &  &  &  &  \\ \hline
                SET 1 &  &  &  &  &  &  \\ \hline
            \end{tabular}
        \end{table}
        上述数据访问一共产生了 \rule{2.5cm}{0.25mm} 次命中。
        \qn 在(2)的基础上增加一条新规则：地址区间包含 5 的倍数的 block 将不会被缓存。仍使用 MRU 替换策略，上述数据访问一共产生了 \rule{2.5cm}{0.25mm} 次命中。
        \qn 在(3)的基础上又增加一条数据预取规则：每当地址为 \verb|10| 的数据被访问时，地址为 \verb|8| 的数据将会被放入缓存，上述数据访问一共产生了 \rule{2.5cm}{0.25mm} 次命中。
        \proy{2015} 某高速缓存满足 $E=2, B=4, S=16$，地址宽度为 14。当引用地址 \verb|0x9D28| 起始的 1 个字节时，tag 位为：
        \begin{choices}
            \item \verb|01110100|
            \item \verb|001110000|
            \item \verb|1110000|
            \item \verb|10011100|
        \end{choices}
        \proy{2015} 关于高速缓存的说法正确的是：
        \begin{choices}
            \item 直写（write through）比写回（write back）在电路实现上更复杂。
            \item 固定的高速缓存大小，较大的块可提高时间局部性好的程序的命中率。
            \item 随着高速缓存组相联度的不断增大，失效率不断下降。
            \item 以上说法全不正确。
        \end{choices}
        \proy{2015} 考虑下面的程序：
        \begin{minted}[frame=single, fontsize=\small]{c}
    #define LENGTH 8
    void clear4x4 (char array[LENGTH][LENGTH]) {
        int row, col;
        for (col = 0; col < 4; col++)
            for (row = 0; row < 4; row++)
                array[row][col] = 0;
    }
        \end{minted}
        \qn 设机器的地址宽度为 7，传入的数组 \verb|array| 的起始地址为 \verb|0x1000000|。两路组相联高速缓存的块大小为 4 字节，一共 4 组，替换算法使用的是 LRU。
            \subqn 以上程序执行会引起多少次失效？
            \subqn 如果 \verb|LENGTH| 改为 16，会引起多少次失效？
            \subqn 如果 \verb|LENGTH| 变为 17，与 b) 相比，下面描述正确的是：\rule{1cm}{0.25mm}，会引起 \rule{2.5cm}{0.25mm} 次失效。
            \begin{choices}
                \item $16 \times 16$ 比 $17 \times 17$ 产生更多的失效次数。
                \item $16 \times 16$ 和 $17 \times 17$ 产生的失效次数相同。
                \item $16 \times 16$ 比 $17 \times 17$ 产生更少的失效次数。
            \end{choices}
            \subqn 请画出 c) 运行后 cache 的最终状态。格式要求：tag 使用二进制；data 使用十进制，例如 \verb|M[6-7]| 表示地址 6 和 7 对应的数据。
            \begin{table}[H]
                \tt
                \centering
                \begin{tabular}{|c|c|c|c|c|c|c|}
                    \hline
                    & {\quad V \quad} & {\quad Tag \quad} & {\quad Data \quad} & {\quad V \quad} & {\quad Tag \quad} & {\quad Data \quad} \\ \hline
                    SET 0 &  &  &  &  &  &  \\ \hline
                    SET 1 &  &  &  &  &  &  \\ \hline
                \end{tabular}
            \end{table}
        \qn 设机器的地址宽度为 8，传入的数组 \verb|array| 的起始地址为 \verb|0x10000000|。全相联高速缓存的块大小为 4 字节，总容量 16 字节，替换算法使用的是 LRU。
            \subqn Tag 字段的位数是 \rule{2.5cm}{0.25mm}。
            \subqn 以上程序执行会引起多少次失效？
            \subqn 如果 \verb|LENGTH| 改为 16，会引起多少次失效？
            \subqn 如果 \verb|LENGTH| 变为 17，与 c) 相比，下面描述正确的是：\rule{1cm}{0.25mm}。
            \begin{choices}
                \item $16 \times 16$ 比 $17 \times 17$ 产生更多的失效次数。
                \item $16 \times 16$ 和 $17 \times 17$ 产生的失效次数相同。
                \item $16 \times 16$ 比 $17 \times 17$ 产生更少的失效次数。
            \end{choices}
            \subqn 请画出 d) 运行后 cache 的最终状态。格式要求：tag 使用二进制；data 使用十进制，例如 \verb|M[6-7]| 表示地址 6 和 7 对应的数据。
            \begin{table}[H]
                \tt
                \centering
                \begin{tabular}{|c|c|c|}
                    \hline
                    V & {\qquad \qquad Tag \qquad \qquad} & {\qquad \qquad Data \qquad \qquad} \\ \hline
                    1 &  &  \\ \hline \hline
                    V & Tag & Data \\ \hline
                    1 &  &  \\ \hline \hline
                    V & Tag & Data \\ \hline
                    1 &  &  \\ \hline \hline
                    V & Tag & Data \\ \hline
                    1 &  &  \\ \hline
                \end{tabular}
            \end{table}
        \proy{2014} 关于 cache 的 miss rate，下面哪些说法是错误的？ 
        \begin{choices}
            \item 保持 $E$ 和 $B$ 不变，增大 $S$，miss rate 一定不会增加。
            \item 保持总容量和 $B$ 不变，增大 $E$，miss rate 一定不会增加。
            \item 保持总容量和 $E$ 不变，增大 $B$，miss rate 一定不会增加。
            \item 如果不采用 LRU，使用随机替换策略，miss rate 可能会降低。
        \end{choices}
        \proy{2014} 现有一个能够存储 4 个 block 的 cache，每一个 cache block 的大小为 \SI{2}{Byte}（$B=2$）。内存空间的大小是 \SI{16}{Byte}，即其地址范围为 \texttt{0\textsubscript{10} (0000\textsubscript{2}) --- 15\textsubscript{10} (1111\textsubscript{2})}。此设定下一程序的访问内存地址序列如下所示，单位是 Byte，数字为十进制：
        \begin{center}
            \tt 2, 3, 10, 9, 6, 8.
        \end{center}
        \qn 若 cache 的结构如下图所示（$S=2, E=2$），初始状态为空，替换策略为 LRU。请在下图空白处填入上述数据访问后 cache 的状态。格式要求：tag 使用二进制；data 使用十进制，例如 \verb|M[6-7]| 表示地址 6 和 7 对应的数据。
        \begin{table}[H]
            \tt
            \centering
            \begin{tabular}{|c|c|c|c|c|c|c|}
                \hline
                & {\quad V \quad} & {\quad Tag \quad} & {\quad Data \quad} & {\quad V \quad} & {\quad Tag \quad} & {\quad Data \quad} \\ \hline
                SET 0 &  &  &  &  &  &  \\ \hline
                SET 1 &  &  &  &  &  &  \\ \hline
            \end{tabular}
        \end{table}
        上述数据访问一共产生了 \rule{2.5cm}{0.25mm} 次 miss。
        \qn 若替换策略为 MRU，其余和(1)相同，请在下图空白处填入上述数据访问后 cache 的状态。
        \begin{table}[H]
            \tt
            \centering
            \begin{tabular}{|c|c|c|c|c|c|c|}
                \hline
                & {\quad V \quad} & {\quad Tag \quad} & {\quad Data \quad} & {\quad V \quad} & {\quad Tag \quad} & {\quad Data \quad} \\ \hline
                SET 0 &  &  &  &  &  &  \\ \hline
                SET 1 &  &  &  &  &  &  \\ \hline
            \end{tabular}
        \end{table}
        上述数据访问一共产生了 \rule{2.5cm}{0.25mm} 次 miss。
        \proy{2013} 如果直接映射高速缓存大小是 \SI{4}{KB}，块大小为 32 字节，则它每组有多少行？
        \begin{choices}
            \item 128
            \item 68
            \item 32
            \item 1
        \end{choices}
        \proy{2013} 现有一个能够存储 4 个 block 的 cache，每一个 cache block 的大小为 \SI{2}{Byte}（$B=2$）。内存空间的大小是 \SI{32}{Byte}，即其地址范围为 \texttt{0\textsubscript{10} (00000\textsubscript{2}) --- 31\textsubscript{10} (11111\textsubscript{2})}。此设定下一程序的访问内存地址序列如下所示，单位是 Byte，数字为十进制：
        \begin{center}
            \tt 1, 4, 17, 2, 8, 16, 9, 0.
        \end{center}
        \qn 如果 cache 的结构是直接映射，满足 $S=4, E=1$，请在下表空白处填入访问上述数据序列访问后 cache 的状态。格式要求：tag 使用二进制；data 使用十进制，例如 \verb|M[6-7]| 表示地址 6 和 7 对应的数据。
        \begin{table}[H]
            \tt
            \centering
            \begin{tabular}{|c|c|c|}
                \hline
                V & {\qquad \qquad Tag \qquad \qquad} & {\qquad \qquad Data \qquad \qquad} \\ \hline
                1 &  &  \\ \hline \hline
                V & Tag & Data \\ \hline
                1 &  &  \\ \hline \hline
                V & Tag & Data \\ \hline
                1 &  &  \\ \hline \hline
                V & Tag & Data \\ \hline
                1 &  &  \\ \hline
            \end{tabular}
        \end{table}
        \qn 如果 cache 的结构如下表所示，满足 $S=2, E=2$，请在下表空白处填入访问上述数据序列访问后 cache 的状态。
        \begin{table}[H]
            \tt
            \centering
            \begin{tabular}{|c|c|c|c|c|c|c|}
                \hline
                & {\quad V \quad} & {\quad Tag \quad} & {\quad Data \quad} & {\quad V \quad} & {\quad Tag \quad} & {\quad Data \quad} \\ \hline
                SET 0 &  &  &  &  &  &  \\ \hline
                SET 1 &  &  &  &  &  &  \\ \hline
            \end{tabular}
        \end{table}
        上述数据访问一共产生了 \rule{2.5cm}{0.25mm} 次 miss。
        \qn 如果 cache 的结构改为 $S=1, E=4$，最终存储在 cache 里面的数据有哪些？只需要填写数据部分，顺序不限。
        \begin{center}
            \rule{2.5cm}{0.25mm}、\quad \rule{2.5cm}{0.25mm}、\quad \rule{2.5cm}{0.25mm}、\quad \rule{2.5cm}{0.25mm}。
        \end{center}
    \end{problems}
	\chapter*{部分参考答案}
    \vspace{-1cm}
    \shipoutAnswer
\end{document}